%% 
%% Copyright 2007-2020 Elsevier Ltd
%% 
%% This file is part of the 'Elsarticle Bundle'.
%% ---------------------------------------------
%% 
%% It may be distributed under the conditions of the LaTeX Project Public
%% License, either version 1.2 of this license or (at your option) any
%% later version.  The latest version of this license is in
%%    http://www.latex-project.org/lppl.txt
%% and version 1.2 or later is part of all distributions of LaTeX
%% version 1999/12/01 or later.
%% 
%% The list of all files belonging to the 'Elsarticle Bundle' is
%% given in the file `manifest.txt'.
%% 

%% Template article for Elsevier's document class `elsarticle'
%% with numbered style bibliographic references
%% SP 2008/03/01
%%
%% 
%%
%% $Id: elsarticle-template-num.tex 190 2020-11-23 11:12:32Z rishi $
%%
%%
\documentclass[review]{elsarticle}

%% Use the option review to obtain double line spacing
%% \documentclass[authoryear,preprint,review,12pt]{elsarticle}

%% Use the options 1p,twocolumn; 3p; 3p,twocolumn; 5p; or 5p,twocolumn
%% for a journal layout:
%% \documentclass[final,1p,times]{elsarticle}
%% \documentclass[final,1p,times,twocolumn]{elsarticle}
%% \documentclass[final,3p,times]{elsarticle}
%% \documentclass[final,3p,times,twocolumn]{elsarticle}
%% \documentclass[final,5p,times]{elsarticle}
%% \documentclass[final,5p,times,twocolumn]{elsarticle}

%% For including figures, graphicx.sty has been loaded in
%% elsarticle.cls. If you prefer to use the old commands
%% please give \usepackage{epsfig}

\usepackage{lineno,hyperref}
\usepackage{amsmath}
\usepackage{amssymb}
\usepackage{mathtools}
\usepackage{mathabx}
\usepackage{hyperref}
\newtheorem{theorem}{Theorem}
\newtheorem{lemma}{Lemma}
\newtheorem{corollary}{Corollary}
\newtheorem{conjecture}{Conjecture}
\newproof{pf}{Proof}
\newdefinition{rmk}{Remark}
\newdefinition{definition}{Definition}
\usepackage{enumitem}
\usepackage{tikz}
\modulolinenumbers[5]

%% The lineno packages adds line numbers. Start line numbering with
%% \begin{linenumbers}, end it with \end{linenumbers}. Or switch it on
%% for the whole article with \linenumbers.
%% \usepackage{lineno}

\journal{---}

\begin{document}

\begin{frontmatter}

%% Title, authors and addresses

%% use the tnoteref command within \title for footnotes;
%% use the tnotetext command for theassociated footnote;
%% use the fnref command within \author or \address for footnotes;
%% use the fntext command for theassociated footnote;
%% use the corref command within \author for corresponding author footnotes;
%% use the cortext command for theassociated footnote;
%% use the ead command for the email address,
%% and the form \ead[url] for the home page:
%% \title{Title\tnoteref{label1}}
%% \tnotetext[label1]{}
%% \author{Name\corref{cor1}\fnref{label2}}
%% \ead{email address}
%% \ead[url]{home page}
%% \fntext[label2]{}
%% \cortext[cor1]{}
%% \affiliation{organization={},
%%             addressline={},
%%             city={},
%%             postcode={},
%%             state={},
%%             country={}}
%% \fntext[label3]{}

\title{Universal Bounds for Spreading on Networks}

%% use optional labels to link authors explicitly to addresses:
%% \author[label1,label2]{}
%% \affiliation[label1]{organization={},
%%             addressline={},
%%             city={},
%%             postcode={},
%%             state={},
%%             country={}}
%%
%% \affiliation[label2]{organization={},
%%             addressline={},
%%             city={},
%%             postcode={},
%%             state={},
%%             country={}}

\author[mymainaddress]{Gadi Fibich\corref{mycorrespondingauthor}}
\cortext[mycorrespondingauthor]{Corresponding author}
\ead{fibich@tau.ac.il}

\author[mymainaddress]{Tomer Levin}


\address[mymainaddress]{Department of Applied Mathematics, School of Mathematical Sciences, Tel Aviv University, Tel Aviv 6997801, Israel}

\begin{abstract}
	Spreading (diffusion) of innovations is a stochastic process on social networks. When the key driving mechanism is peer effects (word of mouth), the rate at which the aggregate adoption level increases with time depends strongly on the network structure. In many applications, however, the  network structure is unknown. To estimate the aggregate adoption level for such innovations, we show that the two networks that correspond to the slowest and fastest adoption levels are a homogeneous two-node network and a homogeneous infinite  complete network, respectively.  Solving the stochastic Bass model on these two networks yields explicit lower and upper bounds for the adoption level on any network. These bounds are tight, and they also hold for the individual adoption probabilities of nodes. The gap between the lower and upper bounds increases monotonically with the ratio of the rates of internal and external influences.
\end{abstract}


\begin{keyword}
Spreading in networks, diffusion of innovations, new products, stochastic models, agent-based models, Bass model
\end{keyword}

\end{frontmatter}

%% \linenumbers
\section{Intorduction}
Diffusion (spreading) in networks is an active research area  
in mathematics, economics, management science, physics, biology, computer science, and social sciences,
and it concerns the spreading of diseases, computer viruses, rumors, information, opinions, technologies,
innovations, and more~\citep{Albert-00,Anderson-92,Jackson-08,Pastor-Satorras-01,Strang-98}. 
In marketing, diffusion of new products is a classical problem~\citep{Mahajan-93}.

The first mathematical model of diffusion of new products
was introduced by Bass~\cite{Bass-69}.  In this model, individuals adopt a new product because
of {\em external influences} by mass media and commercials, and because of {\em internal influences} (peer effect, word-of-mouth)
by individuals who have already adopted the product.
Let~$f$ denote the {\em adoption level} (fraction of adopters) in the population at time~$t$. Then according to the Bass model, %$f$ satisfies the differential equation
	%\begin{equation}
	%\label{eq:homogeneous_Bass}
	%n'(t)= \left(M-n\right)\left(p+\frac{q}{M}n\right), \qquad n(0)=0.
	%\end{equation}
	%In terms of the fraction of adopters, equation~\eqref{eq:homogeneous_Bass} reads
	\begin{equation}
		\label{eq:homogeneous_Bass_frac}
		f'(t)= \left(1-f\right)\left(p+qf\right), \quad t>0, \qquad f(0)=0.
	\end{equation}
	Thus, the $1-f$ potential adopters adopt due to external influences at the constant rate of~$p$, and due to internal influences at the rate of~$qf$, which is proportional to the fraction of adopters.
Equation~\eqref{eq:homogeneous_Bass_frac} can be solved explicitly, yielding the S-shaped Bass formula~\citep{Bass-69}
	\begin{equation}
		\label{eq:Bass_sol}
		f_{\rm Bass}(t)=\frac{1-e^{-(p+q)t}}{1+\frac{q}{p}e^{-(p+q)t}}.
	\end{equation}


The Bass model~\eqref{eq:homogeneous_Bass_frac} inspired a
huge body of theoretical and empirical research; in 2004 it  
was selected as one of the 10 most-cited papers in the 50-year
history of Management Science~\citep{Hopp-04}. Initially, this research was carried out using {\em compartmental Bass models}, such as~\eqref{eq:homogeneous_Bass_frac}, in which the population is divided into several  compartments (e.g., nonadopters and adopters), and the transition
rates between compartments are given by deterministic ordinary differential equations.  {\it Compartmental Bass models, therefore, 
	implicitly assume  that the underlying social network is a homogeneous complete graph}, i.e., that
all individuals within the population are equally likely to influence each other.

In order not to make these assumptions, in more recent studies
diffusion of new products has been studied 
using  Bass models on networks, for the stochastic adoption decision of each individual~\citep{Delre-07, OR-10, Garber-04, Garcia-05, Muller-Peres-19, Rand-11}. 
These {\em agent-based models} allow for implementing a network structure, so that individuals are only influenced by adopters who are also their peers. For example, it has been suggested that social networks have a small-worlds~\citep{watts1998collective} or a scale-free structure~\citep{barabasi1999emergence}. 
In large-scale online social networks,   40\% of the links
were found to be within a 100 km radius~\citep{scellato2011socio}.
In the diffusion of solar panels, the key predictor of a new solar installation is having a neighbor who already installed one~\citep{Bollinger-12,Graziano-15}, and so  the relevant network is a two-dimensional Cartesian grid.  Bass models on networks also allow for heterogeneity among individuals~\cite{PNAS-12,fibich2022diffusion, GLM-01}.


%A mathematical analysis of the  Bass model can be challenging. Hence, most studies that used a  Bass model with a spatial structure were numerical.
   The  effects of 
	{\em various network  characteristics}  (average degree, clustering, \dots)  
on the diffusion were studied numerically using agent-based simulations, see e.g.,~\cite{Goldenberg-10,Goldenberg-08,Peres-14}. For example, it was found that growth is particularly effective in networks that demonstrate cohesion (strong mutual influence among its members), connectedness (high number of ties),
and conciseness (low redundancy) \cite{Muller-Peres-19}.

{\it Explicit expressions} for the expected adoption level~$f(t)$ in the  Bass model were only obtained for a few networks.
Niu \cite{Niu-02} explicitly computed the expected adoption level~$f_{\rm complete}(t;M)$ on complete homogeneous networks with $M$~nodes, and showed that
$\lim_{M \to \infty} f_{\rm complete}(t;M) = f_{\rm Bass}(t)$, see  Theorem~\ref{thrm:Niu} below. 
Fibich and Gibori \cite{OR-10} explicitly computed the expected adoption level~$f_{\rm circle}(t;M)$ on  homogeneous circles with $M$~nodes. They showed that the adoption level on the infinite circle,
 denoted by $f_{\rm 1D}(t):=\lim_{M \to \infty} f_{\rm circle}(t;M)$,
  is given by $f_{\rm 1D}(t)=1-e^{-(p+q)t+q\frac{1-e^{-pt}}{p}}$.

 
  
For most networks,  explicit expressions for~$f(t)$ are not available. Moreover,
in many applications, the network structure or even its characteristics are not known. Hence, {\em it is important, for both theoretical and practical considerations, to obtain explicit lower and upper bounds for the expected adoption level~$f(t)$.} 
%Thus, the adoption level lies between these bounds for any network.  
%In addition, the size of the gap between these bounds  
%indicates the degree to which the network structure can influence the adoption level.


In \cite{OR-10} it was conjectured that since circular and complete networks are the ``least-connected" and the ``most-connected" networks, the adoption level on any infinite network should be bounded from below by that on the infinite circle, and from above by that on the infinite complete network, i.e., that
$
f_{\rm 1D}(t)\leq f(t) \leq f_{\rm Bass}(t). 
$
So far, this conjecture has remained open.


In this study, we settle this conjecture. We 
prove that $f(t)\leq f_{\rm Bass}(t)$ for any finite or infinite network.
Thus, as was conjectured in~\cite{OR-10},  $f_{\rm Bass}$ is a universal upper bound for the adoption level.
Moreover, this upper bound is tight, and is strict for non-complete networks. The tight universal upper bound for 
the individual adoption probabilities of nodes
(i.e., for the probability of any node to adopt the product before time~$t$)
is also given by~$f_{\rm Bass}$.
 
The universal lower bound for~$f(t)$ on general finite or infinite networks, however, is not~$f_{\rm 1D}$. Rather, we prove that $f(t)\geq f_{M=2}^{\rm hom}$ 
for any network,
where~$f^{\rm hom}_{M=2}:= 1-e^{-pt}\frac{qe^{-pt}-pe^{-qt}}{q-p}$ is the expected adoption level on a homogeneous two-node network. 
%\,\footnote{See eq.~\eqref{eq:N_M=2} below.} 
This universal lower bound is also tight, and it also holds for the individual adoption probabilities of nodes. Thus, the conjecture from~\cite{OR-10} that~$f_{\rm 1D}$ is a universal lower bound for all infinite networks is wrong
(note, however, that for any~$D \ge 1$, $f_{\rm 1D}$ is the tight lower bound for the adoption level~$f_{\rm D}(t)$ on infinite D-dimensional Cartesian network 
where each node is connected to its 2D nearest neighbors with edges of weight $\frac{q}{2D}$, 
see~\cite{Funnel2022} for more details).

Let us motivate the ``success" of the conjecture from~\cite{OR-10} regarding the upper bound, and its ``failure" regarding the lower bound. As noted, the compartmental Bass model~\eqref{eq:homogeneous_Bass_frac} corresponds to a complete network, which is indeed the ``most-connected'' network, in the sense that each node can be  directly influenced by all other nodes. A one-sided circle, where each node can only influenced by the node to its left, however, is not the ``least-connected" network. This is because each node is also indirectly influenced by all other nodes. Rather, the ``least-connected" network is a  collection of disjoint pairs of nodes, where each node can be directly influenced by the other node in the pair, but cannot be indirectly influenced by any other node.


To quantify the influence of the social-network structure on the adoption level of new products, we study the {\em size of the gap between the lower and upper bounds}.
The gap size is a monotonically-increasing function of the 
ratio~$\frac{q}{p}$ of the rates of internal and external influences. 
For products that spread predominantly through word of mouth,
we obtain an explicit approximation for the gap size. This explicit approximation shows that the network structure indeed has a large influence on the adoption level of such products. 

   
%In~\citep{Bass-boundary-18}, we obtained an explicit but cumbersome expression for
%the adoption probability of nodes on the finite two-sided line. We then showed by a rather technical proof that in~1D, one-sided diffusion is strictly slower than two-sided diffusion. Using the tools developed in this paper, we obtained a considerably simpler proof of this result, and a much shorter expression for the adoption probability on a finite line. Moreover, we also generalized this result to the anisotropic case, and proved that for multi-dimensional periodic Cartesian grids, the lower bound for the adoption curve is given by that on the circle. 
The paper is organized as follows. Section~\ref{sec:discrete_Bass} presents the   Bass model on a general network. Section~\ref{sec:main_results} presents the main results of this paper on the universal  lower and upper bounds.
Section~\ref{sec:gap} considers the size of the gap between the lower and upper bounds.
Section~\ref{sec:open} lists some open research problems. 
 The detailed proofs are given in Section~\ref{sec:proofs}.


%From a methodological perspective, we also developed the integral Bass inequality, which is used in the proofs of the lower bounds. 
%In addition, we also showed that for networks that have a positive fraction of nodes with a finite indegree, such as finite networks, infinite Cartesian networks, infinite  scale-free network, infinite  small-worlds networks, or infinite Erd\"os-R\'enyi networks, the upper bound is strict. As for the diffusion on a one-dimensional network, it does serve a lower bound for all D-dimensional Cartesian networks. 



%In this study, we provide lower and upper bounds for the diffusion rate on any spatial network, so that the diffusion rate can be at least bounded,  
%In addition, that study conjectured that
%and conjectured that the diffusion on any spatial network is bounded from below and from above by that on a one-dimensional and complete networks, respectively. 
%and in~\cite{Bass-SIR-analysis-17} for the discrete Bass-SIR model.



%Additional studies 
%--- The effect of boundary
%conditions on the diffusion in Cartesian networks was studied in~\citep{Bass-boundary-22, Bass-boundary-18}. Fibich
%and Levin~\citep{fibich2020percolation} analyzed the percolation of new products on various networks, from which a fraction
%of the nodes is randomly removed. The effect of consumers heterogeneity in various spatial networks was studied in~\citep{fibich2022diffusion}.




%Our motivation comes from 
%products that diffuse predominantly through internal influences by geographical neighbors,
%such as residential rooftop solar systems~\cite{Bollinger-12,Graziano-15}. For such products, it is reasonable to approximate the social network by a two-dimensional grid, in which each node represents a residential unit. Previous studies of the  Bass model 
%on two-dimensional networks avoided boundary effects by 
%imposing periodic boundary conditions~\cite{Bass-SIR-model-16,Bass-SIR-analysis-17,OR-10}. 
%In that case, all nodes are interior and are influenced  from all sides by the adjacent nodes. 
%In practice, however, some residential units (nodes) lie at the boundary of the municipality. In addition, some 
%interior residential units  are separated by a physical barrier (river, highway) from adjacent units. 
%Therefore, real two-dimensional networks are not periodic, but rather have exterior and possibly also interior boundaries.}


\section{Bass model on networks}
	\label{sec:discrete_Bass}
	
	We begin by introducing the  Bass model on a general heterogeneous network. A new product is introduced at time $t=0$ to a network with 
	$M$~individuals, denoted by~${\cal M}:=\{1,\ldots,M\}$, where $M$ can be finite or infinite. We denote by~$X_j(t)$ the {\em state} of individual $j$ at time $t$, so that 
	\begin{equation*}
		X_j(t)=\begin{cases}
			1, \quad {\rm if}\ j\ {\rm is\ an\ adopter\ of\ the\ product\ at\ time}\ t,\\
			0, \quad {\rm otherwise.}
		\end{cases}
	\end{equation*} 
	Since the product is new, all individuals are initially nonadopters, i.e., 
	\begin{subequations}
	\label{eq:dbm}
	\begin{equation}
		\label{eq:general_initial}
		X_j(0)=0, \qquad j\in {\cal M}.
	\end{equation}
 The underlying social network is represented by a {\em weighted directed} graph, such that if there is an  edge from $k$ to $j$, the rate of {\em internal influence} of adopter~$k$ on nonadopter~$j$ to adopt is~$q_{k,j}>0$, and~$q_{k,j}=0$ if there is no edge from~$k$ to~$j$.
 	The edges and influence rates are {\em not} assumed to be symmetric, i.e., 
 $q_{k,j}$~may be different from~$q_{j,k}$. 
Since nonadopters do not self-influence to adopt,
$$
q_{j,j}\equiv 0, \qquad j\in {\cal M}.
$$	
In contrast to similar models in epidemiology on networks~\citep{kiss2017mathematics}, such as the Susceptible Infected (SI) model, $j$ also experiences {\em external influences}  to adopt by mass media and commercials, at a constant rate of~$p_j>0$. 
Internal and external influences are assumed to be additive. Thus,  
{\em the adoption time~$T_j$ of nonadopter~$j$ is exponentially distributed} at the rate of
\begin{equation}
\label{eq:lambda_j}
\lambda_j(t):=p_j+\sum_{k\in {\cal M}}q_{k,j}X_k(t), \qquad j\in {\cal M}, \quad t>0,
\end{equation}
which increases whenever $k$~adopts and $q_{k,j}>0$. Finally, it is assumed that once an individual adopts the product, she or he remains an adopter for all time.
Therefore, the stochastic adoption of~$j\in {\cal M}$ in the time interval~$(t,t+\Delta t)$ as~$\Delta t\rightarrow 0$ is given by
	\begin{equation}
		\label{eq:general_model}
		\mathbb{P}(X_j(t+\Delta t)=1\,| \, {\bf X}(t))=
		\begin{cases}
			\mbox{} \qquad \qquad  1,\hfill &\quad {\rm if}\ X_j(t)=1,\\
			\left(p_j+\sum\limits_{\substack{k \in {\cal M}}}q_{k,j}X_{k}(t)\right)\Delta t, &\quad {\rm if}\ X_j(t)=0,
		\end{cases}
	\end{equation}
\end{subequations}
	where ${\bf X}(t) := \left(X_1(t),\ldots, X_M(t)\right)$ is the {\em state of the network} at time~$t$. Note that the time variable is continuous.

	
	
	The maximal rate of internal influences that can be exerted on node~$j$ (which is when all its neighbors/peers  are adopters) is 
	\begin{subequations}
\begin{equation}
		\label{eq:q_on_node}
		q_j:=\sum_{\substack{k \in {\cal M}}}q_{k,j}.
	\end{equation}
	For simplicity, we assume that each node can be influenced by at least one node, i.e.,
	\begin{equation}
	\label{eq:q_j>0}
	q_j>0, \qquad j\in {\cal M}.
	\end{equation}
\end{subequations}	
 We do not assume, however, that the network  only consists of  a single connected component. 
 	The underlying network of the  Bass model~\eqref{eq:dbm} is denoted by
	\begin{equation}
		\label{eq:def-cal-N}
	{\cal N}={\cal N}({\cal M}, \{p_k\}_{k\in {\cal M}}, \{q_{k,j}\}_{k,j\in {\cal M}}).
	\end{equation}


The adoption level at time~$t$ is~$\frac{1}{M}\sum_{\substack{j \in {\cal M}}}X_j(t)$. Our goal is to 
obtain lower and upper bounds for
the {\bf expected adoption level} (fraction of adopters)
	\begin{equation*}
		\label{eq:number_to_fraction}
		f(t;{\cal N}):=\frac{1}{M}\mathbb{E} \Big[\sum_{\substack{j \in {\cal M}}}X_j(t)\Big].
	\end{equation*}
To do that, we will compute lower and upper bounds for  the adoption probabilities of nodes
\begin{equation*}
f_j(t;{\cal N}) := \mathbb{P}(X_j(t)=1)=\mathbb{E}\left[X_j(t)\right], \qquad j\in {\cal M},
\end{equation*}
and then use
\begin{equation}
	\label{eq:f=sum-f_j}
f= \begin{cases}
	\frac{1}{M}\sum_{j=1}^M f_j, & \quad M<\infty, \\
	\lim\limits_{M\rightarrow \infty}\frac{1}{M}\sum_{j=1}^M f_j, & \quad M = \infty.
\end{cases}
\end{equation}
The dependence of the adoption level and of the adoption probabilities of nodes on the external and internal influence rates is monotonic:
\begin{theorem}[\cite{Bass-boundary-18}]
	\label{thm:dominance-principle-f_j}
	Consider the  Bass model~\eqref{eq:dbm} on network~$\cal N$, see~\eqref{eq:def-cal-N}. Let $t>0$. Then~$f(t;{\cal N})$ is monotonically increasing, and~$\{f_m(t;{\cal N})\}$ are monotonically non-decreasing, with respect to each~$p_j$ and each~$q_{k,j}$.
\end{theorem}	
%Therefore, adding edges increases increases its weight from zero to a positive weight,
%\begin{corollary}
%	\label{cor:dominance-principle-f}
%	If network~${\cal N}_B$ is obtained from network~${\cal N}_B$ by adding edges, then $f(t;{\cal N}_B)>f(t;{\cal N}_A)$ for $0<t<\infty$.
%\end{corollary} 

\subsection{Homogeneous complete networks}

Let $f_{\rm complete}(t;p,q,M)$ denote the expected adoption level  in the  Bass model~\eqref{eq:dbm} on the homogeneous complete network ${\cal N}_{\rm complete}(p,q,M)$, defined as 
\begin{equation}
	\label{eq:comp_net}
	p_j \equiv  p, \quad q_{k,j}=\begin{cases}
		\frac{q}{M-1}, \qquad &k\neq j,\\
		0, \qquad & k =j,
	\end{cases}, \qquad j,k\in {\cal M}.
\end{equation}
As $M$ increases, each node is influenced by more nodes, 
but the weight of each node decreases, so that the maximal rate of internal influences $q_j \equiv q$  remains unchanged, see~\eqref{eq:q_on_node}. 
Nevertheless, the expected adoption level increases with~$M$: 
\begin{lemma}[\cite{Bass-monotone-convergence-23}]
	\label{lem:f_complete-monotone-in-M}
	Let $t,p,q>0$. Then $f_{\rm complete}(t;p,q,M)$ is monotonically increasing 
	in~$M$.
\end{lemma}
As~$ M \to \infty$, the  Bass model~\eqref{eq:dbm} on complete networks
approaches the original compartmental Bass model:
%Let us recall the following result, which is due to Niu:
\begin{theorem}[\cite{Niu-02}]
	\label{thrm:Niu}
	$
	\lim\limits_{M\rightarrow \infty} f_{\rm complete}(t;p,q,M) = f_{\rm Bass}(t;p,q),
	$
	where~$f_{\rm Bass}$ is given by~\eqref{eq:Bass_sol}.
\end{theorem}
     From Lemma~\ref{lem:f_complete-monotone-in-M} and Theorem~\ref{thrm:Niu} we have
\begin{corollary}
	\label{cor:f_complete<f_Bass}
	Let $t,p,q>0$. Then 
	\begin{equation*}
		\label{eq:f_complete<f_Bass}
		f_{\rm complete}(t;p,q,M)< f_{\rm Bass}(t;p,q), \qquad  M = 1,2, \dots
	\end{equation*} 
\end{corollary}
%Therefore, 
%if $f \le f_{\rm complete}$, then $f$ is bounded from above by $f_{\rm Bass}$,
%as was conjectured in~\citep{OR-10}.


%\subsection{Monotonicity in $p_j$ and in $q_{k,j}$}


%
%
%Denote 
\section{Main results}
\label{sec:main_results}
In this section we present the main results of this paper. The proofs are given in Section~\ref{sec:proofs}.
For clarity, we formulate the results for networks  that are homogeneous in~$\{p_j\}$ and~$\{q_j\}$, i.e.,
\begin{equation}
	\label{eq:p_j,q_j-homogeneous}
	p_j\equiv p, \quad q_j \equiv q,\qquad j\in{\cal M}.
\end{equation}
This requirement can be satisfied by any graph structure that satisfies~\eqref{eq:q_j>0}, 
and not just by the complete network~\eqref{eq:comp_net}. For example, for any given network ${\cal N}({\cal M},\{p_j\}, \{q_{k,j}\})$, define network~$\widetilde{{\cal N}}({\cal M},\{\widetilde{p_j}\}, \{\widetilde{q}_{k,j}\})$ such that~$\widetilde{p_j}:=p$ and $\widetilde{q}_{k,j}:=q_{k,j}\frac{q}{q_j}$. Then~$\widetilde{{\cal N}}$ satisfies~\eqref{eq:p_j,q_j-homogeneous}, and it has the same nodes/edges structure as~$\cal N$.

The extension of the results to networks which do not satisfy~\eqref{eq:p_j,q_j-homogeneous} is discussed in Section~\ref{sec:inhomogeneous-pj-qj}. We also note that, quite often,
the difference in~$f$ between a network which is heterogeneous in~$\{p_j\}$ and~$\{q_j\}$
and the corresponding network which is homogeneous in~$\{p_j\}$ and~$\{q_j\}$
is quite small, even when the level of heterogeneity is not~\citep{OR-10,PNAS-12}.


\subsection{Non-tight universal bounds}

The following  universal lower and upper bounds are immediate:
\begin{lemma}
	  \label{lem:trivial-bounds}
		Consider the Bass model~\eqref{eq:dbm}
	on a network~${\cal N}$
%	$
%	{\cal N} = {\cal N}({\cal M}, \{p_k\}_{k\in {\cal M}}, \{q_{k,j}\}_{k,j\in {\cal M}})
%	$ 
	which is homogeneous in~$\{p_j\}$ and~$\{q_j\}$, see~\eqref{eq:p_j,q_j-homogeneous}. 
	Then 
	\begin{subequations}
\begin{equation}
   \label{eq:trivial-bounds-f_m}
	1-e^{-p t} \le f_m(t) \le 1-e^{-(p+q)t}, 
	 \quad t \ge 0, \quad m \in {\cal M},
\end{equation}
and so
\begin{equation}
	\label{eq:trivial-bounds-f}
1-e^{-p t} \le f(t) \le 1-e^{-(p+q)t}, 
    \quad  t \ge 0.
\end{equation}
	\end{subequations}		
\end{lemma}
\begin{pf}
Since $X_k(t) \in \{0,1 \}$ for any~$k \in \cal M$,  the adoption rate of node~$m$ is bounded by, see~\eqref{eq:lambda_j} and~\eqref{eq:q_on_node},
$$
p=p_m \le \lambda_m(t) \le p_m+\sum_{k\in {\cal M}}q_{k,m} = p+q, \quad m\in {\cal M}, \quad t \ge 0.
$$
Hence, \eqref{eq:trivial-bounds-f_m}~follows, and so 
\eqref{eq:trivial-bounds-f}~follows by~\eqref{eq:f=sum-f_j}.
%
%	$$
%1-e^{-p_j t} \le f_j(t) \le 1-e^{-(p_j+q_j)t}, \qquad j\in {\cal M}, \quad t>0.
%$$
\end{pf}

Thus, the lower and upper bounds~\eqref{eq:trivial-bounds-f_m}  for~$f_m(t)$  
correspond to 
%are tight, we note that they  
%also follow from Theorem~\ref{thm:dominance-principle-f_j}.
%Indeed, the lower bound in~\eqref{eq:trivial-bounds-f_m} correspond
%to letting  $q_m \searrow 0$, and the upper bound correspond
%to letting  $p_j \to \infty$ for all $j \not= m$ (so that all nodes but $m$ adopts at
%$t=0+$). 
the extreme cases when none of the other individuals adopted by time $t$,
and when all the other individuals adopted at~$t = 0+$, respectively. 
Therefore, these bounds are not expected to be tight, as indeed we will show below. 
 



\subsection{Tight upper bound}
\label{sec:upper-bound}


If one adds edges to a network, this increases the adoption level~$f$ (Theorem~\ref{thm:dominance-principle-f_j}). 
The following two observations suggest a stronger result, namely, that
even if as we add edges, we lower the weights of the edges so as to keep $q_j \equiv q$  unchanged, the adoption level increases:
 \begin{enumerate}
 	\item The adoption level $f_{\rm complete}(t;p,q,M)$ in homogeneous complete networks is monotonically increasing in~$M$ (Lemma~\ref{lem:f_complete-monotone-in-M}).
 	
 	\item  The adoption level~$f_D(t;p,q)$ in infinite $D$-dimensional Cartesian networks, where each node is connected to its 
 	$2D$ nearest neighbors, and the weights of these edges  is~$\frac{q}{2D}$,
 	is monotonically increasing in~$D$ (this was shown numerically and asymptotically in~\cite{OR-10}).
 	  
 \end{enumerate}
 Thus,    %for networks  that satisfy~\eqref{eq:p_j,q_j-homogeneous}, 
{\em numerous weak edges lead to a faster diffusion than a few strong ones}.
 Therefore, we can expect  that among all networks with $M$~nodes that satisfy~\eqref{eq:p_j,q_j-homogeneous}, the fastest diffusion would be on the complete network ${\cal N}_{\rm complete}(p,q,M)$, see~\eqref{eq:comp_net}, as formulated in Conjecture~\ref{conj-ub-M} below. If that is indeed the case,  then by Corollary~\ref{cor:f_complete<f_Bass}, the adoption levels 
on all networks should be bounded from above by~$f_{\rm Bass}$.
  %
%
%
%Consider the Bass model~\eqref{eq:dbm} on a network~${\cal N}$
%which is homogeneous in~$\{p_j\}$ and~$\{q_j\}$, see~\eqref{eq:p_j,q_j-homogeneous}
%By Theorem~\ref{thm:dominance-principle-f_j} and Corollary~\ref{cor:dominance-principle-f}, 
%the adoption level on~${\cal N}$ is bounded from above by
%that on a complete network with the same number of nodes, the same weights for the original edges, and any positive weight for the new edges. 
% That upper bound, however, can be quite crude. Intuitively, this is because   
% simply adding edges increases the values of~$\{q_j\}$ above~$q$. 
%The results of~\cite{OR-10} suggest, however, that adding edges while keeping the
%$\{q_j\}$ unchanged (by lowering the weights of the original edges)
%leads to an increase in the adoption level.
%In other words, {\em numerous weak internal influences 
%	lead to  faster diffusion than a few strong ones}. 
%Therefore, we formulate a new conjecture: 
%\begin{conj}
%	\label{conj-ub-M}
%	The maximal adoption level among all networks with $M$~nodes which are homogeneous in~$\{p_j\}$ and~$\{q_j\}$, see~\eqref{eq:p_j,q_j-homogeneous}, is attained  on the complete network ${\cal N}_{\rm complete}(p,q,M)$, i.e.,  
%	$$
%	\max_{\{{\cal N} \, | \, \text{${\cal N}$ has $M$ nodes and satisfies \eqref{eq:p_j,q_j-homogeneous}} \}} f(t;{\cal N})  =   f_{\rm complete}(t;p,q,M), \qquad t \ge 0.
%	$$ 	
%\end{conj} 	
%
%%Let $${\cal G}(p,q,M):=\{{\cal N} \, | \, \text{${\cal N}$ has $M$ nodes, \eqref{eq:p_j,q_j-homogeneous} holds} \}$$
%%be the set of
%%
%%  In the informal opening of Section~\ref{sec:upper-bound}, we conjectured that the fastest diffusion in~${\cal G}(p,q,M)$ occurs on the homogeneous complete network~\eqref{eq:comp_net}:
%
%
%% 
%%  To illustrate this, consider a homogeneous one-sided circle with $M$ nodes, where each node can be
%% influenced by its left neighbor at the rate of~$q$, and  a two-sided circle with $M$ nodes, where each node can be influenced by its left and right neighbors at the rates of $q_L$ and~$q_R$, respectively. If $q = q_L+q_R$, then 
%% the adoption levels on the one-sided and two-sided circles are identical, i.e.,  
%% $f_{\rm circle}^{\rm one-sided}(t;p,q) = f_{\rm circle}^{\rm two-sided}(t;p,q_L,q_R,M)$,
%% see~\cite{OR-10}. Therefore, a more meaningful comparison of networks is 
%% to keep $q_j$, see~\eqref{eq:q_on_node}.
%% For example,  
%%by $f^D$ the adoption level in the Bass model on a $D$-dimensional Cartesian
%%%grid, such that each node is connected to its 2D adjacent nodes, and all edges 
%%%have the same weight of~$q^D$. Then   $f^D$ is bounded by 
%%Let us set the 
%%weight of each edge to $\frac{q}{2D}$. Then as the dimension $D$ increases,
%%each node is influenced by mode nodes, but the weight of each edge decreases,
%%so that $q_j$ remains unchanged. In this case, it was shown numerically and asymptotically
%%that $f^D$ increases with~$D$, see~\cite{OR-10}. 
%Note that it provides a better (i.e. lower) upper bound than the one obtained by simply adding edges, since the weights of the complete network are lower than 
%in the proginla bound.  
%%This suggests that the fastest adoption would be on a complete network with $q_j \equiv q$, which is a tighter 
%%estimate than a complete network where the weight of each node is~$q$.
%
% 
%    
%
%
%Consider again the Bass model~\eqref{eq:dbm} on network~${\cal N}$ with $M$~nodes, which is homogeneous in~$\{p_j\}$ and~$\{q_j\}$.
%%\begin{equation}
%%\label{eq:p_j,q_j-homogeneous}
%%p_j\equiv p, \quad q_j \equiv q,\qquad j\in{\cal M}.
%%\end{equation}
%If Conjecture~\ref{conj-ub-M} holds, then 
%% the  adoption level~$f(t;{\cal N})$  would be maximal when the influence~$q$ on any node in the network is equally split among all other~$M-1$ nodes.
%%If this assumption is correct, then for a given~$M$, the fastest diffusion is on the complete network~\eqref{eq:comp_net}.
%%%where
%%% \begin{equation}
%%% \label{eq:comp_net}
%%% p_j \equiv  p, \quad q_{k,j}=\begin{cases}
%%%			\frac{q}{M-1}, \qquad &k\neq j,\\
%%%			0, \qquad & k =j,
%%%		\end{cases}, \qquad j,k\in {\cal M}.
%%% \end{equation} 
%% Thus, we assume that 
% $f(t;{\cal N}) \le f_{\rm complete}(t;p,q,M)< f_{\rm Bass}(t;p,q)$, 
%%  where
%% $f_{\rm complete}(t;p,q,M)$ is the expected adoption level in the  Bass model~\eqref{eq:dbm} on the  complete network~\eqref{eq:p_j,q_j-homogeneous}. 
% where the second inequality follows from Corollary~\ref{cor:f_complete<f_Bass}
%   Therefore, 
%  $f$ is bounded from above by $f_{\rm Bass}$,
%  as was conjectured in~\citep{OR-10}.
%  
%%  Hence, the diffusion should be bounded by that on an infinite complete network, which is given by~$f_{\rm Bass}$, see Theorem~\ref{thrm:Niu} below.
%Of course, proving one conjecture by formulating a new one does not really 
%make a difference. We now prove the old conjecture that, indeed,  
%the adoption probabilities of nodes are indeed bounded from above by~$f_{\rm Bass}$. Hence, ,  
%the expected adoption level is bounded from above by~$f_{\rm Bass}$
%without making use of the new conjecture:
Indeed, we can rigorously prove 
\begin{theorem}
	\label{thm:bass_largest}
	Consider the Bass model~\eqref{eq:dbm}
	on a network~${\cal N}$
%	$
%	{\cal N} = {\cal N}({\cal M}, \{p_k\}_{k\in {\cal M}}, \{q_{k,j}\}_{k,j\in {\cal M}})
%	$ 
	which is homogeneous in~$\{p_j\}$ and~$\{q_j\}$, see~\eqref{eq:p_j,q_j-homogeneous}. 
Then
	\begin{equation}
		\label{eq:f_m<=f_Bass}
		f_{m}(t;{\cal N}) \leq f_{\rm Bass}(t;p,q), \qquad t \ge 0, \qquad m \in {\cal M},
	\end{equation}
	where~$f_{\rm Bass}$ is given by~\eqref{eq:Bass_sol}, and so
	\begin{equation}
		\label{eq:f<=f_Bass}
		f(t;{\cal N})\leq f_{\rm Bass}(t;p,q), \qquad t \ge 0.
	\end{equation}
\end{theorem}


 In Lemma~\ref{lem:trivial-bounds} we derived the upper bound
$
f_j(t), f(t) \le 1- e^{-(p+q)t}$.
The upper bound
%$f_j(t), f(t) \le  f_{\rm Bass}(t; p,q)$
of Theorem~\ref{thm:bass_largest} is better (i.e., lower), since by~\eqref{eq:Bass_sol},
$$
f_{\rm Bass}(t;p,q)=\frac{1-e^{-(p+q)t}}{1+\frac{q}{p}e^{-(p+q)t}}< 1- e^{-(p+q)t}.
$$ 
We can further show that $f_{\rm Bass}$ is the {\em tight} universal upper bound:
%In light of Theorems~\ref{thrm:Niu} and~\ref{thm:bass_largest}, $f_{\rm Bass}$ is the {\em tight} universal upper bound, i.e., there is no better (i.e. smaller) upper bound which holds for all networks that satisfy 
%$p_j \leq  p$ and $q_{j} \leq q$ for $j \in {\cal M}$:
\begin{lemma}
	\label{lem:f_Bass-is-tight}
	The universal upper bound in Theorem~\ref{thm:bass_largest} is tight, in the sense that
	% 	$$
	% 	   \sup_{\{{\cal N}({\cal M},\{p_j\},\{q_{k,j}\}) \, | \, \text{\eqref{eq:p_j,q_j<p,q} holds}\}, m \in {\cal M}} f_m(t;\{p_j\},\{q_{k,j}\}) = f_{\rm Bass}(t;p,q),
	% 	$$
	\begin{eqnarray*}
	&\sup_{\{ \, {\cal N} \, | \, \text{\eqref{eq:p_j,q_j-homogeneous} holds}\}} f(t;{\cal N})
	 =
	  %f_{\rm Bass}(t;p,q),  \qquad &
	  \sup_{\{ \, {\cal N} \, | \, \text{\eqref{eq:p_j,q_j-homogeneous} holds}\}, \, m \in \cal M} f_m(t;{\cal N}) = f_{\rm Bass}(t;p,q).
	\end{eqnarray*}	
%	where~${\cal G}^{U}(p,q):=\{ \, {\cal N} \, | \, \text{\eqref{eq:p_j,q_j<p,q} holds}\}$ is the set of all networks that satisfy~\eqref{eq:p_j,q_j<p,q}.
\end{lemma}
  While the upper bound~$f_{\rm Bass}$ is attained for an infinite homogeneous complete network (Theorem~\ref{thrm:Niu}), it is strict for nodes that have a finite indegree, hence for networks with a positive fraction of nodes with finite indegree:
 \begin{theorem}
	\label{thm:f<f_Bass}
Assume the conditions of Theorem~\ref{thm:bass_largest}. 
\begin{enumerate}
\item   If node~$m$ has a finite indegree,
	then 
	\begin{equation}
		\label{eq:f_m<f_Bass}
		f_m(t;{\cal N}) < f_{\rm Bass}(t;p,q), \qquad t>0.
	\end{equation}

\item If there is a positive fraction of nodes  in the network with a finite indegree, then
\begin{equation}
	 \label{eq:f<f_Bass=thm}
	f(t;{\cal N}) < f_{\rm Bass}(t;p,q), \qquad t > 0.
\end{equation}
\end{enumerate}
\end{theorem}
  
  
  
Therefore, the upper bound~$f_{\rm Bass}$ is strict for any network which is not  infinite and complete (up to a vanishing fraction of nodes). In particular, assume that the network type is one of the following: 
\begin{itemize}
	\item A finite network.
	\item An infinite (homogeneous or heterogeneous) D-dimensional Cartesian network.
	\item An infinite  scale-free network~\citep{barabasi1999emergence}.
	\item An infinite  small-worlds network~\citep{watts1998collective}.
	\item The infinite sparse random networks $\lim_{M \to \infty} G\left(M,\frac{\lambda}{M}\right)$ ~\citep{erdHos1960evolution}.
\end{itemize}
	Since all these finite and infinite networks have a positive fraction of finite-indegree nodes, Theorem~\ref{thm:f<f_Bass} implies that $f<f_{\rm Bass}$ for all these networks types. 


\subsection{Tight lower bound}

Let ${\cal N}_{M=2}^{\rm hom}(p,q)$ denote the
 homogeneous network with two nodes, where
\begin{equation}
\label{eq:N_M=2}
{\cal M} = \{1, 2\}, \quad  p_1=p_2 = p, \quad q_{1,2} = q_{2,1} = q, \quad q_{1,1} = q_{2,2} = 0.
\end{equation}
The expected adoption level on ${\cal N}_{M=2}^{\rm hom}$ can be explicitly calculated  (see, e.g.,~\cite{Bass-boundary-18}),  giving
\begin{equation}
\label{eq:f_all-M=2}
f_{M=2}^{\rm hom}(t;p,q)  = 1-e^{-pt}\frac{qe^{-pt}-pe^{-qt}}{q-p}, \qquad p \neq q.
\end{equation}
Note that there is only one homogeneous network with two nodes. Thus,
$f_{M=2}^{\rm hom}(t;p,q)  = f_{\rm complete}^{\rm hom}(t;p,q,M=2)  = f_{\rm circle}^{\rm hom}(t;p,q,M=2)$.

The informal arguments at the beginning of Section~\ref{sec:upper-bound} 
  suggest that {\em few strong edges lead to a slower diffusion than numerous   weak ones}. Hence,  
it is intuitive to expect that for given~$p$ and~$q$, the adoption level is lowest when the influence~$q$ on any node in the network is exerted by a single node.  
 %  There are, however, numerous networks in which each node is only influenced by a single node at the rate of~$q$.    
   	This requirement is satisfied when the network is a one-sided circle, or a collection of disjoint one-sided circles. Among all circles, the lowest adoption is on a two-node circle (Lemma~\ref{lem:f_complete-monotone-in-M}). Intuitively, this is because 
   on a two-node circle each node can only be influenced by one node, whereas on longer circles each node can also be indirectly influenced by additional nodes. 
%   In what follows,
%   we  prove that the minimal adoption level is indeed attained on the two-node circle.
%   
%   
%   Two extremes cases are a one-sided circle, where each node is also indirectly influenced by all other nodes, and a collection of two-node networks of type~${\cal N}_{M=2}^{\rm hom}$, in which there are no such indirect influences. Therefore, the diffusion in ${\cal N}_{M=2}^{\rm hom}$ should be the slowest among all networks in which each node is influenced by a single node, hence among all networks.
%
%
%
%
Indeed, we now prove that $f_{M=2}^{\rm hom}$ is a universal lower bound for~$\{f_m\}$, hence for~$f$:
\begin{theorem}
\label{thm:f_m>=f_M=2}
Assume the conditions of Theorem~\ref{thm:bass_largest}. 
Then
\begin{subequations}
	\label{eqs:f,f_m>=f_M=2}
	\begin{equation}
		\label{eq:f_m>=f_M=2}
		f_m(t;{\cal N})  \ge  f_{M=2}^{\rm hom}(t;p,q), \qquad t \ge 0, \qquad m \in {\cal M},
	\end{equation}
	and so
	\begin{equation}
		\label{eq:f>=ff_M=2}
		f(t;{\cal N})  \ge  f_{M=2}^{\rm hom}(t;p,q), \qquad t \ge 0.
	\end{equation}
\end{subequations}
\end{theorem}

In Lemma~\ref{lem:trivial-bounds}, we derived the lower bound
$
f_j(t), f(t)  \ge  1-e^{-pt}$. The lower bound
in Theorem~\ref{thm:f_m>=f_M=2} is better (i.e., larger), since by Theorem~\ref{thm:dominance-principle-f_j},
$$
f_{M=2}^{\rm hom}(t;p,q)>f_{M=2}^{\rm hom}(t;p,q=0) = 1-e^{-pt}.
$$ 

Moreover, $f_{M=2}^{\rm hom}$ is the {\em tight} universal lower bound:
\begin{lemma}
	\label{lem:f_M=2-is-tight}
	Let $M \in\{ 2,4, \dots \}$. 
	Then
	% 	$$
	% 	   \sup_{\{{\cal N}({\cal M},\{p_j\},\{q_{k,j}\}) \, | \, \text{\eqref{eq:p_j,q_j<p,q} holds}\}, m \in {\cal M}} f_m(t;\{p_j\},\{q_{k,j}\}) = f_{\rm Bass}(t;p,q),
	% 	$$
	\begin{eqnarray*}
&\inf_{\{ \, {\cal N} \, | \, \text{\eqref{eq:p_j,q_j-homogeneous} holds}\}} f(t;{\cal N}) = %f_{M=2}^{\rm hom}(t;p,q),  \qquad   &
  \inf_{\{ \, {\cal N} \, | \, \text{\eqref{eq:p_j,q_j-homogeneous} holds}\}, \, m \in \cal M} f_m(t;{\cal N}) = f_{M=2}^{\rm hom}(t;p,q).	
	\end{eqnarray*}
%where~${\cal G}^{L}(p,q):=\{ \, {\cal N}  \, | \, \text{\eqref{eq:p_j,q_j>=p,q} holds}\}$ is the set of all networks that satisfy~\eqref{eq:p_j,q_j>=p,q}.
\end{lemma}


 The lower bound~$f,f_j\ge  f_{M=2}^{\rm hom}$ is attained 
  if the network is a collection of disjoint pairs of nodes, 
  each of which is of type~${\cal N}_{M=2}^{\rm hom}$. For all other networks,  however,  it is strict:
\begin{theorem}
	\label{thm:f_j>f_M=2^hom}
	Assume the conditions of Theorem~\ref{thm:bass_largest}. 
	\begin{itemize}
		\item 
		If node~$j$ belongs to a connected component with more than two nodes, then 
		\begin{equation}
			\label{eq:f_m>f_2nodes}
			f_j(t;{\cal N}) > f_{M=2}^{\rm hom}(t;p,q), \qquad t>0.
		\end{equation}
		
		\item 	If the fraction of nodes  in~$\cal N$ that belong to a connected component with more than two nodes is positive, then
	\begin{equation}
		\label{eq:f>f_M=2^hom-thm}
		f(t;{\cal N})> f_{M=2}^{\rm hom}(t;p,q), \qquad t > 0.
	\end{equation}
			\end{itemize}
\end{theorem}




\subsection{Bounds for networks   inhomogeneous in 
	$\{p_j\}$ or~$\{q_j\}$}
  \label{sec:inhomogeneous-pj-qj}

We can extend all the upper-bound results to networks which are not homogeneous in 
$\{p_j\}$ and in~$\{q_j\}$, as follows:
\begin{corollary}
	Theorem~\ref{thm:bass_largest}, Lemma~\ref{lem:f_Bass-is-tight}, 
	and  Theorem~\ref{thm:f<f_Bass},  
	remain valid if we replace  condition~\eqref{eq:p_j,q_j-homogeneous} with 
	\begin{equation}
		\label{eq:p_j,q_j<p,q}
		p_j \leq  p, \quad q_j \leq q, \qquad j\in {\cal M}.
	\end{equation} 
\end{corollary} 
\begin{pf}
This follows from Theorem~\ref{thm:dominance-principle-f_j}.
%	By the dominance principle,
%	$\{f_m\}$ and~$f$ are monotonically increasing with respect to each~$p_j$ 
%	and each~$q_{k,j}$  (Lemma~\ref{lem:monotone_in_p_and_q}). 
%	Therefore, any upper bound obtained for~\eqref{eq:p_j,q_j-homogeneous}, remains valid
%	if~\eqref{eq:p_j,q_j<p,q} holds.   
	%Hence, it is sufficient to show that~\eqref{eq:f_m<=f_Bass} 
	%and~\eqref{eq:f<=f_Bass} hold in the homogeneous case~\eqref{eq:p_j,q_j-homogeneous}.
\end{pf}

Similarly, we can extend all the lower-bound results to networks which are not homogeneous in 
$\{p_j\}$ and in $\{q_j\}$: 
\begin{corollary}
	Theorem~\ref{thm:f_m>=f_M=2}, Lemma~\ref{lem:f_M=2-is-tight}, and  Theorem~\ref{thm:f_j>f_M=2^hom}, 
	remain valid if we replace condition~\eqref{eq:p_j,q_j-homogeneous} with 
	\begin{equation}
		\label{eq:p_j,q_j>=p,q}
		p_j \geq  p, \quad q_j \geq q, \qquad j\in {\cal M}.
	\end{equation}
\end{corollary}


\section{Gap between lower and upper bounds}
\label{sec:gap}

Consider any network~${\cal N}$ which is homogeneous in~$\{p_j\}$ and~$\{q_j\}$, see~\eqref{eq:p_j,q_j-homogeneous}.
By Theorems~\ref{thm:bass_largest} and~\ref{thm:f_m>=f_M=2}, the expected adoption level and the adoption probability of nodes are bounded by
  $$
f^{\rm hom}_{M=2}(t;p,q)\leq f(t;{\cal N}), f_m(t;{\cal N}) \leq f_{\rm Bass}(t;p,q),
\qquad t\ge 0.
$$  
Therefore, it is natural to consider the {\em size of the gap} between the explicit lower and upper bounds~$f_{M=2}^{\rm hom}$ and~$f_{\rm Bass}$,
which expresses the dependence of the diffusion on the network structure. 


The explicit bounds can be written in a dimensionless form as
  $$
f_{M=2}^{\rm hom}(t;p,q)=f_{M=2}^{\rm hom}\left(\widetilde{t};\widetilde{q}\right), \quad f_{\rm Bass}(t;p,q) = f_{\rm Bass}\left(\widetilde{t};\widetilde{q}\right),
$$
where~$\widetilde{t}=qt$ and~$\widetilde{q}=\frac{q}{p}$.
The nondimensional parameter~$\widetilde{q}$  expresses the ratio of internal and external influences. 
Since network effects are only due to internal influences, they increase with~$\frac{q}{p}$. Thus,
when~$q=0$, there are no network effects, and so the two bounds are identical, i.e., 
$$
f^{\rm hom}_{M=2}(t;p,q=0) = f_{\rm Bass}(t;p,q=0)=1-e^{-pt}.
$$
 When~$\frac{q}{p}\ll 1$,
 the network has a minor effect on the diffusion, and so~$f^{\rm hom}_{M=2}\approx f_{\rm Bass}$, see Figure~\ref{fig:f_bass_2nodes}A.  For products that spread predominantly through word-of-mouth, however, the regime of relevance is~$\frac{q}{p} \gg 1$, typically $10\leq \frac{q}{p}\leq 100$~\citep{Bass-69}. As can be expected, the difference between~$f^{\rm hom}_{M=2}$ and $f_{\rm Bass}$ is significant for~$\frac{q}{p}=10$ (Figure~\ref{fig:f_bass_2nodes}B), and even larger for~$\frac{q}{p}=100$ (Figure~\ref{fig:f_bass_2nodes}C). Note that for any network~${\cal N}$, $f(t;{\cal N})$ lies in the shaded region between~$f^{\rm hom}_{M=2}(t)$ and $f_{\rm Bass}(t)$. 

\begin{figure*}[ht!]
\begin{center}
\scalebox{1}{\includegraphics{f_bass_2nodes.eps}}
\end{center}
\caption{The expected adoption level~$f(t)$ of any network~${\cal N}$ lies in the shaded region between the lower bound~$f^{\rm hom}_{M=2}$ (orange dash) and upper bound~$f_{\rm Bass}$ (blue solid). 
 (A)~$\frac{q}{p}=0.1$. The two bounds are nearly indistinguishable. (B)~$\frac{q}{p}=10$. (C)~$\frac{q}{p}=100$. (D)~$\frac{T_{\rm Bass}^{1/2}}{T^{1/2, {\rm hom}}_{M=2}}$ as a function of~$\frac{q}{p}$ (blue solid), and its asymptotic approximation~\eqref{eq:approx} (black dashes). }
\label{fig:f_bass_2nodes}
\end{figure*}

It is  instructive to compare the adoption levels on different networks using the ``half-life''~$T^{1/2}$ for half of the population to adopt.
In particular, we can use~$T^{1/2}$ to compare the bounds~$f_{\rm Bass}$ and~$f_{M=2}^{\rm hom}$.
 The ratio~$\frac{T_{\rm Bass}^{1/2}}{T^{1/2, {\rm hom}}_{M=2}}$ can be estimated asymptotically, yielding
 \begin{equation}
 \label{eq:approx}
 \frac{T_{\rm Bass}^{1/2}}{T^{1/2, {\rm hom}}_{M=2}}\sim \frac{2}{\log 2 }\frac{p}{q}\log \frac{q}{p}, \qquad  \frac{q}{p} \gg 1.
 \end{equation}
 Figure~\ref{fig:f_bass_2nodes}D confirms that~$\frac{T^{1/2}_{\rm Bass}}{T^{1/2, {\rm hom}}_{M=2}}$ decreases with~$\frac{q}{p}$, and approaches the asymptotic limit~\eqref{eq:approx} as~$\frac{q}{p}\rightarrow \infty$.
 %Thus, $\frac{T^{1/2}_{\rm Bass}}{T^{1/2, {\rm hom}}_{M=2}}\approx 0.9992$ for~$\frac{q}{p}=0.1$, $\frac{T^{1/2}_{\rm Bass}}{T^{1/2, {\rm hom}}_{M=2}}\approx 0.6$ for~$\frac{q}{p}=10$,
%and~$\frac{T^{1/2}_{\rm Bass}}{T^{1/2, {\rm hom}}_{M=2}}\approx 0.1$ for~$\frac{q}{p}=100$.
%\section{Discussion}
%\label{sec:discussion}
%The results follow from the following ``principles'':
%\begin{enumerate}
%\item For given network type, $f$ increases with~$M$.
%\item In general, adoption in homogeneous networks is faster than in heterogeneous netowrks.
%\item Having~$q$ fixed, more peers leads to faster adoption.
%\end{enumerate}
%Let us consider a network with~$M$ nodes
% that is homogeneous in~$\{p_j\}$ and~$\{q_j\}$, see~\eqref{eq:p_j,q_j-homogeneous}. In general, heterogeneity slows down the diffusion, see~\citep{fibich2022diffusion}. Therefore, the fastest diffusion will be when a node is influenced by all the~$M-1$ nodes at the rate of~$\frac{q}{M-1}$, and the slowest when it is influenced by just on node, i.e., $q_{k,j}=q\delta_{k,i}$, for some~$i\ne j$. Hence,  $f(t;p,q_{k,j}=q\delta_{k,i})\leq f(t;p,q)\leq f_{\rm complete}(t;p,q,k)$. In addition, having more peers (with a fixed~$q$) leads to faster diffusion. Therefore, $f_{\rm complete}(t;p,q, M)$ is monotonically increasing in~$M$, see~\citep{COMPARTMENTAL2022}.
%Therefore, the upper bound is~$\lim_{M\rightarrow \infty}f_{\rm complete} = f_{\rm Bass}$. As for the lower bound, if~$j$ is influenced just by~$i$, than the slowest adoption would occur if~$i$ is not influenced by anyone but~$j$. Hence, we get~$f_{M=2}^{\rm hom}$.
This limit goes to zero as~$\frac{q}{p}\rightarrow \infty$, showing that the network structure has a large effect on the diffusion when~$\frac{q}{p}\gg 1$, i.e., for products that  diffuse primarily by internal influences.


%Since~$f^{\rm hom}_{M=2}(t;p,q) = f_{\rm circle}(t;p,q,M=2)$, $f_{\rm 1D}(t;p,q) = \lim_{M\rightarrow \infty} f_{\rm circle}(t;p,q,M)$, and since~$f_{\rm circle}$ is monotonically increasing in~$M$, see~\citep{Funnel2022}, then~$f^{\rm hom}_{M=2}(t;p,q)<f_{\rm 1D}(t;p,q)$. Therefore, $

\section{Open problems}
  \label{sec:open}
  
This manuscript settles the conjecture from~\cite{OR-10},
 but  leads to some new questions, which are currently open. 
 Indeed, the upper and lower bounds in Theorems~\ref{thm:bass_largest}
 and~\ref{thm:f_m>=f_M=2} are tight for networks with any number of nodes. 
Can these bounds be improved if we restrict ourselves to networks with a fixed number of nodes? 

Thus, let $${\cal G}(p,q,M):=\{{\cal N}  \mid \text{${\cal N}$ has $M$ nodes, \eqref{eq:p_j,q_j-homogeneous} holds} \}$$
 be the set of all networks with $M$~nodes that are homogeneous in~$\{p_j\}$ and~$\{q_j\}$.  In the beginning  of Section~\ref{sec:upper-bound}, we argued that the fastest diffusion in~${\cal G}(p,q,M)$ should occur on the homogeneous complete network~\eqref{eq:comp_net}. Therefore, we formulate 
 \begin{conjecture}
 	\label{conj-ub-M}
 $
\sup_{{\cal N} \in {\cal G}(p,q,M)} f(t;{\cal N})  =   f_{\rm complete}(t;p,q,M).
$ 	
\end{conjecture} 	
%One could try to prove this result as in Thereom~\ref{thm:bass_largest},
%by showing that $[S] \ge  [S_{\rm complete}](t;p,q,M)$. The difficulty here is that while $S_{\rm Bass}$ is a solution of a single ODE, $[S_{\rm complete}](t;p,q,M)$ is the solution of a system of ODES.
We note, however, that the rate of convergence of~$f_{\rm complete}$ to~$f_{\rm Bass}$  as $M \to \infty$ is~$O(\frac1M)$, see~\cite{Bass-compartmental-limits-22}. Therefore, the difference between these two upper bounds becomes negligible for large (e.g., $M = 10^6$) networks.
 
 
 Consider now the lower bound.
 Let $M$ be even, and let network~$\tilde{\cal N}$ be composed of~$\frac{M}{2}$ pairs of nodes, each of which is of type~${\cal N}_{M=2}^{\rm hom}$, see~\eqref{eq:N_M=2}. Then $f(t;\tilde{\cal N}) = f_{M=2}^{\rm hom}(p,q)$. 
Therefore,
$$
\inf_{{\cal N} \in {\cal G}(p,q,M)} f(t;{\cal N})  =  f_{M=2}^{\rm hom}(p,q), 
\qquad \text{$M$ even}.
$$
Thus,  the lower bound $f_{M=2}^{\rm hom}$ cannot be improved 
(i.e., increased) for networks with a fixed even number of nodes.
The tight lower bound for $M$~odd, however, is an open problem. 

 Another open question is the tight lower bound of~$f$ among {\em connected}
networks with $M$~nodes (even or odd) that are homogeneous in~$\{p_j\}$ and~$\{q_j\}$.
Here one may need to distinguish between connected undirected networks, 
 weakly-connected directed graphs (there is an undirected path between any pair of vertices), and strongly-connected directed graphs (there is a directed path between every pair of vertices).


% Is it given by~$f_{\rm circle}(t;p,q,M)$?


\section{Proof of results}
\label{sec:proofs}


\subsection{Master equations}

Denote the nonadoption probability of node~$j$ by
\begin{equation}
\label{eq:nonadoption_prob}
[S_j](t):=1-f_j(t)= \mathbb{P}(X_j(t)=0).
\end{equation}
Then $[S_j]$ satisfies
 the {\em master equation}~\citep{fibich2022diffusion}
	\begin{equation}
		\label{eq:master_Sa}
		\frac{d}{dt}[S_j](t)= -\left(p_j+q_j\right)\left[S_{j}\right]
		+\sum_{\substack{k \in {\cal M}}}q_{k,j}[S_j,S_k](t), \qquad [S_{j}](0) =1,
	\end{equation} 
where~$q_j$ is given by~\eqref{eq:q_on_node}, and
$$
[S_j,S_k](t):=\mathbb{P}(X_j(t)=X_k(t)=0).
$$
In general, to close these equations, one adds the master equations for all pairs~$\{[S_j, S_k]\}$, all triplets~$\{[S_j, S_k, S_m]\}$, etc., see~\cite{fibich2022diffusion}. For the purpose of obtaining the lower and upper bounds, however, we will only need the following result:
\begin{lemma}
Consider the  Bass model~\eqref{eq:dbm}. Then for any~$i,j \in {\cal M}$,
\begin{equation}
\label{eq:set_e-2pt}
[S_i](t)\,[S_j](t) \leq [S_i,S_j](t) \leq e^{-2pt}, \qquad 0\leq t<\infty.
\end{equation}
\end{lemma}
\begin{pf}
The left inequality is proved in~\cite{Funnel2022}. For the right inequality, we note that the joint nonadoption probability of a pair~$\{i,j\}$ of isolated nodes ($q_{j} = q_{i} = 0$) is given by
		\begin{equation*}
			[S_i,S_j](t;p,q_{j} = q_{i} = 0) = e^{-2pt}, 
		\end{equation*}
see, e.g.,~\cite{Bass-boundary-18}. Hence, the right inequality follows from Theorem~\ref{thm:dominance-principle-f_j}.
\end{pf}
\subsection{Differential and integral Bass inequalities}
Let us recall the following result:
\begin{lemma}[\cite{OR-10}]
	\label{lem:Bass-inequality}
	Let~$p,q>0$, and let $f(t)$ satisfy the {\bf differential Bass inequality}
	$$
	\frac{df}{dt} < (1-f)(p+qf), \quad t>0, \qquad f(0)=0. 
	$$
	Then $f(t) < f_{\rm Bass}(t;p,q)$ for $0<t<\infty$.
\end{lemma}

Let $[S_{\rm Bass}]:=1-f_{\rm Bass}$ denote the nonadoption level in the compartmental
Bass model. Then by~\eqref{eq:homogeneous_Bass_frac},
\begin{equation}
	\label{eq:homogeneous_Bass_frac-[S]}
	\frac{d}{dt}[S_{\rm Bass}](t)%= -[S_{\rm Bass}]\left(p+q(1-[S_{\rm Bass}])\right)
	 =-(p+q)[S_{\rm Bass}]+ q[S_{\rm Bass}]^2  , \quad [S_{\rm Bass}](0)=1.
\end{equation}
If we replace the equality sign in~\eqref{eq:homogeneous_Bass_frac-[S]} by an inequality, the solution of this inequality is bounded from below by~$[S_{\rm Bass}]$:
 \begin{lemma}
 \label{cor:diffBass}
Let~$p,q>0$, and let $[S](t)$ satisfies the {\bf differential Bass inequality}
\begin{equation*}
	\label{eq:homogeneous_Bass_frac-[S]-inequality}
	\frac{d}{dt}[S](t)%= -[S_{\rm Bass}]\left(p+q(1-[S_{\rm Bass}])\right)
	>-(p+q)[S]+ q[S]^2, \quad t>0, \qquad [S](0)=1.
\end{equation*}
 Then $[S](t)> [S_{\rm Bass}](t)$ for $0<t<\infty$.
 \end{lemma}
\begin{pf}
This follows from Lemma~\ref{lem:Bass-inequality}  and~$[S_{\rm Bass}]=1-f_{\rm Bass}$.
\end{pf}



   
Multiplying~\eqref{eq:homogeneous_Bass_frac-[S]} by $e^{(p+q)t}$, integrating between zero and~$t$, and using the initial condition, gives the {\em integral form of the compartmental Bass model}
\begin{equation}
	\label{eq:homogeneous_Bass_frac-S-integral}
	[S_{\rm Bass}] (t) = e^{-(p+q)t} + q\int_0^t e^{-(p+q)(t-\tau)} [S_{\rm Bass}]^2(\tau) \, d\tau.
\end{equation}
If we replace the
equality sign in~\eqref{eq:homogeneous_Bass_frac-S-integral} by an inequality, the solution of the resulting  integral Bass inequality is bounded from below by~$[S_{\rm Bass}]$:
\begin{lemma}
	\label{lem:Bass-inequality-integral}
	Let $p,q>0$, and let $[S](t)$ be non-negative and continuous in~$[0, \infty)$. 
	\begin{enumerate}
		\item 
	If $[S]$ satisfies the {\bf  integral Bass inequality}
\begin{equation}
	\label{eq:S-Bass-inequality-integral}
	[S] (t) \ge  e^{-(p+q)t} + q\int_0^t e^{-(p+q)(t-\tau)} [S]^2(\tau) \, d\tau, \quad t>0,
%\quad \text{{\bf (integral Bass  inequality)}}
\end{equation}
	then 
	$
	[S] (t) \ge 	[S_{\rm Bass}] (t;p,q)$ for $t \ge 0$.

	
	\item If  inequality~\eqref{eq:S-Bass-inequality-integral} is strict, then  
	$[S] (t) > 	[S_{\rm Bass}] (t;p,q)$ for $t>0$.
		\end{enumerate}
\end{lemma}
\begin{pf}
Let $u:=[S]-[S_{\rm Bass}]$. Subtracting~\eqref{eq:homogeneous_Bass_frac-S-integral} from~\eqref{eq:S-Bass-inequality-integral} gives
$$
  u(t) \ge q \int_0^t e^{-(p+q)(t-\tau)}  \left([S]^2-[S_{\rm Bass}]^2\right)(\tau) \, d\tau.
$$
Therefore,
\begin{eqnarray}
	\label{eq:u>=int-phi*u}
	u(t) \ge \int_0^t\phi(\tau) u(\tau) \, d\tau, 
	%\\&
	 \qquad \phi(\tau):=q e^{-(p+q)(t-\tau)}  \left([S]+[S_{\rm Bass}]\right)(\tau).%\nonumber
\end{eqnarray}
%
%
%  where
%  $$
%  \phi(\tau):=q e^{-(p+q)(t-\tau)}  \left([S](\tau)+[S_{\rm Bass}](\tau)\right).
%$$
Since $[S]$ and $[S_{\rm Bass}]$ are continuous and non-negative, then so is~$\phi$.
Let 
\begin{equation}
	\label{eq:v-def-integral-Bass}
 v(t):=e^{-\int_0^t \phi} \int_0^t\phi(\tau) u(\tau) \, d\tau.
\end{equation}
 Then
$v(0) = 0$  and
$$
  \frac{dv}{dt} = 
  e^{-\int_0^t \phi} \phi(t) 
  \left(u(t)- \int_0^t\phi(\tau) u(\tau) \, d\tau \right) \ge 0,
$$
where the inequality follows from~\eqref{eq:u>=int-phi*u}. Therefore,  for $t \ge 0$,  $v(t) \ge 0$.
Hence, by~\eqref{eq:v-def-integral-Bass}, $\int_0^t\phi(\tau) u(\tau) \, d\tau \ge 0$ and so by~\eqref{eq:u>=int-phi*u}, $u(t) \ge 0$.

If inequality~\eqref{eq:S-Bass-inequality-integral} is strict, we replace in the above proof all
 ``$\ge$'' signs by~``$>$'' signs.
\end{pf}



	
\subsection{Upper bound}
 \label{sec:upper}
 

%{\bf Proof of Lemma~\ref{cor:f_complete<f_Bass}. }
%Since $f_{\rm complete}(t;p,q,M)$ is strictly monotonically increasing in~$M$,
%see~\cite{COMPARTMENTAL2022},
%the result follows from Theorem~\ref{thrm:Niu}. \qed

 
 
%To prove that~$\{f_m\}_{m\in{\cal M}}$ and~$f$ are bounded from above by~$f_{\rm Bass}$, 
%In particular, if the network is 
%homogeneous in $p_j \equiv p$ and $q_j^{\rm in} \equiv q$, then $\{f_k\}_{k=1}^M$ and~$f$ are bounded from 
%above by~$f_{\rm Bass}(t;p,q)$. 
%Consider the discrete heterogeneous Bass model~\eqref{eq:Bass-model-heterogeneous-tools}.
%Let $f_k =f_k(t;\{p_j\},\{q_{k,j}\})$ and~$f =\frac1M \sum_{k \in \cal M} f_k$ denote the adoption probability of node~$k$ and the expected adoption level, respectively.
%In this section we use the master equations and the integral Bass inequality to show that if
%\begin{equation}
%	\label{eq:p_j,q_j<p,q}
%	p_j \leq  p, \quad q_j \leq q, \qquad j\in {\cal M},
%\end{equation} 
%where $q_j:=\sum\limits_{k \in \cal M} q_{k,j}$,
%then
%$\{f_k\}_{k \in \cal M}$ and~$f$ are bounded from above by $f_{\rm Bass}(t;p,q)$. 
%In particular, if the network is 
%homogeneous in $p_j \equiv p$ and $q_j \equiv q$, then $\{f_k\}_{k \in \cal M}$ and~$f$ are bounded from 
%above by~$f_{\rm Bass}(t;p,q)$. 
%To prove this upper bound, we  show that 
%any function that satisfies the {\em integral Bass inequality}, see~\eqref{eq:S-Bass-inequality-integral}, is bounded from below by~$[S_{\rm Bass}]:=1-f_{\rm Bass}$. 
%Then we show that 
%$[S_k]:=1-f_k$ satisfies the  integral Bass inequality
%for any $k \in {\cal M}$. 
%Therefore, all  $\{[S_k] \}_{k \in \cal M}$ are bounded from below by~$[S_{\rm Bass}]$, hence also $%[S]:=
%\frac1M \sum_{k \in \cal M} [S_k] = 1-f$.  
We begin with an auxiliary result.
\begin{lemma}
	\label{cor:lipschitz}
	Consider the  Bass model~\eqref{eq:dbm}.
	Let~\eqref{eq:p_j,q_j<p,q} hold, and let
	\begin{equation}
	\label{eq:s_inf_s_j}
		[\underbar{S}](t) :=  \inf_{j\in {\cal M}} \left\{[S_j](t) \right\}.
	\end{equation}
	Then $[\underbar{S}](t)$ is non-negative and continuous.
\end{lemma}
\begin{pf}
	The non-negativity of~$[\underbar{S}]$ follows from that of $\{[S_j]\}$. 
	Let $j \in {\cal M}$.
	Since all probabilities are bounded between~$0$ and~$1$, then using~\eqref{eq:master_Sa} and~\eqref{eq:p_j,q_j<p,q},
	\begin{equation*}
		\left|\frac{d}{dt}[S_j]\right| \leq (p+q)[S_j]+\sum_{k\in \cal M}q_{k,j}\left[S_j,S_k\right]\leq p+q+\sum_{k\in \cal M}q_{k,j} \le  \kappa, 
	\end{equation*}
	%Since boundedness of the first derivative implies Lipschitz continuity, it follows that the functions~$\left\{[S_i]\right\}_{i\geq 1}$ are uniformly Lipschitz continuous with a joint Lipschitz constant~$C:=p+2q$.
	where~$\kappa:=p+2q$.
	Hence, by the mean-value theorem, for any $t,t^*>0$, $\left|[S_j](t) - [S_j]\left(t^*\right)\right|\leq \kappa\left|t-t^*\right|$, and so~$- [S_j]\left(t^*\right)\leq -[S_j](t)+\kappa\left|t-t^*\right|\leq -[\underbar{S}](t)+\kappa\left|t-t^*\right|$. Taking the supremum of the left-hand side yields $ -[\underbar{S}]\left(t^*\right)\leq -[\underbar{S}](t)+\kappa\left|t-t^*\right|$, and so~$[\underbar{S}](t) -[\underbar{S}]\left(t^*\right)\leq \kappa\left|t-t^*\right|$. Swapping~$t$ and~$t^*$ gives the inverse estimate, and so~$[\underbar{S}](t)$ is continuous.
\end{pf}

{\bf Proof of Theorem~\ref{thm:bass_largest}. }
%	By Theorem~\ref{thm:dominance-principle-f_j}, it is sufficient to show that~\eqref{eq:f_m<=f_Bass} 
%	and~\eqref{eq:f<=f_Bass} hold in the homogeneous case~\eqref{eq:p_j,q_j-homogeneous}.
%	In addition, 
Since~$1-f_m = [S_m]\geq [\underbar{S}]$, 
see~\eqref{eq:nonadoption_prob} and~\eqref{eq:s_inf_s_j}, 	
	it is sufficient to show that
	\begin{equation}
		\label{eq:s_bass<s_inf}
		[\underbar{S}](t) \geq [S_{\rm Bass}](t,p,q).
	\end{equation}
	
	By~\eqref{eq:master_Sa} with $q_j = q$, see~\eqref{eq:p_j,q_j-homogeneous}, 
	\begin{equation}
		\label{eq:S_j-imtegral-identity-upper-bound}
		[S_j] = e^{-(p+q)t}+\int_0^te^{-(p+q)\left(t-\tau\right)}
		\sum_{k\in \cal M} q_{k,j} \, [S_j,S_k](\tau)\, d\tau .
	\end{equation}
	Therefore,  by the lower bound in~\eqref{eq:set_e-2pt} and~\eqref{eq:s_inf_s_j},
	\begin{eqnarray*}
			&[S_j] \geq e^{-(p+q)t}+\int_0^te^{-(p+q)(t-\tau)}\sum\limits_{k\in \cal M}q_{k,j} \, [S_j](\tau) \, [S_k](\tau) \, d\tau \\
			&\geq e^{-(p+q)t}+q \int_0^te^{-(p+q)(t-\tau)} [\underbar{S}]^2(\tau) d\tau.
	\end{eqnarray*}
	Taking the infimum over all~$j$ gives
	\begin{equation*}
		%\label{eq:S_inf_integral}
		[\underbar{S}] \geq e^{-(p+q)t}+q \int_0^t e^{-(p+q)(t-\tau)}[\underbar{S}]^2(\tau)\, d\tau.
	\end{equation*}
	Therefore, since $[\underbar{S}]$ is non-negative and continuous (Lemma~\ref{cor:lipschitz}),
	we can use the integral Bass inequality (Lemma~\ref{lem:Bass-inequality-integral})
	to get inequality~\eqref{eq:s_bass<s_inf}, from which~\eqref{eq:f_m<=f_Bass} follows. Therefore, by~\eqref{eq:f=sum-f_j}, 
%	Finally, since $f = \frac{1}{M}\sum_{j=1}^M f_j$ for finite networks and~$f = \lim_{M\rightarrow \infty }\frac{1}{M}\sum_{j=1}^M f_j$ for infinite networks, 
\eqref{eq:f<=f_Bass}~follows.~\qed


{\bf Proof of Lemma~\ref{lem:f_Bass-is-tight}. }
The result for~$f$ follows from  Theorem~\ref{thrm:Niu}.  Since the complete network~\eqref{eq:comp_net} is homogeneous, 
	$f_m \equiv f$ for all $m \in \cal M$. Hence, the result holds for any~$f_m$ as well. \qed




{\bf Proof of Theorem~\ref{thm:f<f_Bass}. }
	Let 
$$
A_{d}({\cal N}):= \{m \in {\cal M} \mid \text{indegree}\left(m\right)=d \}
$$
denote  the set of all nodes with
indegree~$d$ in network~${\cal N}$. Then it is sufficient to prove that for all networks 
that satisfy~\eqref{eq:p_j,q_j-homogeneous}  and for all $d \in \mathbb{N}$, 
\begin{equation}
	\label{eq:S_m>S_Bass-induction}
	[S_m](t;{\cal N}) > [S_{\rm Bass}](t;p,q), \quad t>0, \qquad m \in A_d({\cal N}).
\end{equation}

We prove~\eqref{eq:S_m>S_Bass-induction} by induction on~$d$. 
When $d=0$, node $m \in A_0$ is not influenced by any other node, and so 
\begin{equation}
	\label{eq:e^-pt>S_Bass}
	[S_m](t;{\cal N})= e^{-pt}   =[S_{\rm Bass}](t;p,q=0)
	> [S_{\rm Bass}](t;p,q),
\end{equation}
where the inequality follows from Theorem~\ref{thm:dominance-principle-f_j}.

For the induction stage, we assume that~\eqref{eq:S_m>S_Bass-induction}
holds  for all networks 
that satisfy~\eqref{eq:p_j,q_j-homogeneous}  and for all~$m \in A_{d-1}$, and prove that it holds 
for all networks that satisfy~\eqref{eq:p_j,q_j-homogeneous}  and for all~$m \in A_{d}$, as follows. 
Let $m \in A_d$, where $d \ge 1$, and denote by 
$\{k_1,\ldots, k_{d}\}$ the $d$~nodes that can influence~$m$. 
The master equation for~$[S_m]$ is,
see~\eqref{eq:p_j,q_j-homogeneous} and~\eqref{eq:master_Sa}, 
% see~\eqref{eq:master_Sa1}, %and~\eqref{eq:p_j,q_j-homogeneous}, 
\begin{equation}
	\label{eq:s^j_evolution_J+1}
	\frac{d}{dt}[S_m]= -(p+q)[S_m]
	+\sum_{i=1}^{d} q_{k_i,m}[S_m,S_{k_i}], \qquad [S_m](0)=1 .
\end{equation}
By the indifference principle, we can compute each of the $d$~probabilities $\{[S_m,S_{k_i}]\}_{i=1}^{d}$  on a modified network~$\widetilde{\cal N}_i$, in which we remove the edge~$k_i \rightarrow {m}$. Thus, $	[S_m,S_{k_i}] = \widetilde{[S_m,S_{k_i}]}$,  where the tilde sign refers to probabilities in~$\widetilde{\cal N}_i$. In this modified network, node~$m$ has indegree~$d-1$, and so by the induction assumption\footnote{In the modified network~$\widetilde{\cal N}_i$ we reduced~$q_m$ by~$q_{k_i,m}>0$. Therefore, $\widetilde{q}_m<q$,
	and so we cannot apply the induction assumption directly for~$\widetilde{\cal N}_i$.  
	By Theorem~\ref{thm:dominance-principle-f_j}, however,
	aince the induction assumption holds when~$\widetilde{q}_m=q$, 
	see~\eqref{eq:p_j,q_j-homogeneous}, it also holds when~$\widetilde{q}_m<q$.}
$$
\widetilde{[{S}_m]} > [S_{\rm Bass}]. 
$$
In addition, by Theorem~\ref{thm:bass_largest}, 
$$
\widetilde{[{S}_{k_i}]}  \ge  [S_{\rm Bass}], 
$$
Combining the above and~\eqref{eq:set_e-2pt},
% using Theorem~\ref{thm:[Si,Sj]>=[Si][Sj]} on the modified network~$\widetilde{\cal N}_i$, 
we have that
$$
[S_m,S_{k_i}] = \widetilde{[S_m,S_{k_i}]} \ge \widetilde{[{S}_m]}\widetilde{[{S}_{k_i}]} >[S_{\rm Bass}]^2.
$$
Therefore, 
\begin{equation}
	\label{eq:s^js^i_indifference_J+1}
	\sum_{i=1}^{d} q_{k_i,m} \, [S_m,S_{k_i}] > \sum_{i=1}^{d} q_{k_i,m} [S_{\rm Bass}]^2 = q [S_{\rm Bass}]^2.
\end{equation}

By~\eqref{eq:s^j_evolution_J+1} and~\eqref{eq:s^js^i_indifference_J+1},
\begin{equation*}
	%\label{eq:u>0strict_J+1}
	\frac{d}{dt}[S_m]+(p+q)[S_m]>q[S_{\rm Bass}]^2, \qquad [S_m](0)=1 .
\end{equation*}
This is the differential Bass inequality (Lemma~\ref{lem:Bass-inequality}), written in terms
of~$[S]$, see~\eqref{eq:homogeneous_Bass_frac-[S]-inequality}.  Hence, 
$[S_m] > [S_{\rm Bass}]$, as needed.
%	By Theorem~\ref{thm:dominance-principle-f_j}, it is sufficient to prove the result  when~\eqref{eq:p_j,q_j-homogeneous} holds.
%	Denote by~$A_{d} \subset \{1, \dots, M\}$ the set of all nodes with
%	indegree equal to~$d$. Then it is sufficient to prove that for all networks 
%	that satisfy~\eqref{eq:p_j,q_j-homogeneous}  and for all $d \in \mathbb{N}$, 
%	\begin{equation}
%		\label{eq:S_m>S_Bass-induction}
%		[S_m](t;\{p_j\},\{q_{k,j}\}) > [S_{\rm Bass}](t;p,q), \quad t>0, \quad m \in A_d.
%	\end{equation}
%	
%	We prove~\eqref{eq:S_m>S_Bass-induction} by induction on the indegree~$d$. 
%When $d=0$, node $m \in A_0$ is not influenced by any other node, and so by Theorem~\ref{thm:dominance-principle-f_j},
%	\begin{equation*}
%		\label{eq:e^-pt>S_Bass}
%	[S_m](t;\{p_j\},\{q_{k,j}\})= e^{-p_m t}= e^{-pt}   %=[S_{\rm Bass}](t;p,q=0)
%	 > [S_{\rm Bass}](t;p,q).
%	\end{equation*}	
%	
%	For the induction stage, we assume that~\eqref{eq:S_m>S_Bass-induction}
%	holds for~$d-1$, and show that it holds for~$d$. Let $m \in A_d$, where $d \ge 1$, and
%	let nodes~$\{k_1,\ldots, k_{d}\}$ be the $d$~neighbors of~$m$. 
%	The master equation for~$[S_m]$ is, see~\eqref{eq:p_j,q_j-homogeneous} and~\eqref{eq:master_Sa}, 
%	\begin{equation}
%		\label{eq:s^j_evolution_J+1}
%		\frac{d}{dt}[S_m]= -(p+q)[S_m]
%		+\sum_{i=1}^{d} q_{k_i,m}[S_m,S_{k_i}], \quad [S_m](0)=1 .
%	\end{equation}
%	By the {\em indifference principle} \citep{Bass-boundary-18}, we can compute each of the $d$~probabilities $\{[S_m,S_{k_i}]\}_{i=1}^{d}$  on a modified network~$\widetilde{\cal N}$ in which we remove the edge~$k_i \rightarrow m$, i.e., $	[S_m,S_{k_i}] = \widetilde{[S_m,{S}_{k_i}]}$,  where the tilde sign refers to probabilities in~$\widetilde{\cal N}$. In~$\widetilde{\cal N}$, node~$m$ has $d-1$~neighbors, and so by the induction assumption\footnote{In the modified network we reduced~$q_m$ by~$q_{k_i,m}>0$, and so~\eqref{eq:p_j,q_j<p,q} holds, rather than~\eqref{eq:p_j,q_j-homogeneous}. 
%		Since the induction assumption holds when~\eqref{eq:p_j,q_j-homogeneous} holds, however, it also holds when~\eqref{eq:p_j,q_j<p,q} holds.}
%	$$
%	\widetilde{[{S}_m]} > [S_{\rm Bass}](t;p,q). 
%	$$
%	
%	By Theorem~\ref{thm:bass_largest}, 
%	$$
%	\widetilde{[{S}_{k_i}]}  \ge  [S_{\rm Bass}](t;p,q), 
%	$$
%	Combining the above and~\eqref{eq:set_e-2pt}, we have that
%	\begin{equation}
%		\label{eq:s^js^i_indifference_J+1}
%		[S_m,S_{k_i}] = \widetilde{[S_m,{S}_{k_i}]} \ge \widetilde{[{S}_m]}\widetilde{[{S}_{k_i}]} >[S_{\rm Bass}]^2(t;p,q).
%	\end{equation}
%	By~\eqref{eq:s^j_evolution_J+1} and~\eqref{eq:s^js^i_indifference_J+1},
%	\begin{equation*}
%		%\label{eq:u>0strict_J+1}
%		\frac{d}{dt}[S_m](t)+(p+q)[S_m]>q[S_{\rm Bass}]^2(t;p,q), \qquad [S_m](0)=1 .
%	\end{equation*}
%	Therefore, the differential Bass inequality (Corollary~\ref{cor:diffBass}) implies that
%	$[S_m] > [S_{\rm Bass}]$, as desired.
%Finally, \eqref{eq:f<f_Bass=thm} follows from Theorem~\ref{thm:bass_largest}, \eqref{eq:f=sum-f_j},  and~\eqref{eq:f_m<f_Bass}. 
\qed




%\begin{proof}
%	Let us choose some $m \in {\cal M}$. Define
%	$$
%	\tilde{p}_j = p, \quad 
%	%j \in {\cal M}, \qquad 
%	\tilde{q}_{k,j} = 
%	\begin{cases}
%		0, & \quad \text{if~} j \not=m, 
%	\\
%	{q}_{k,j}& \quad \text{if~} j =m, 
%\end{cases}
%    \qquad k,j \in {\cal M}.
%	$$
%	By the dominance principle for nodes, $f_m(t) \ge \tilde{f}_m(t)$.
%	In Lemma~\ref{lem:aux-f_m>=f_M=2} below we will prove that 
%\begin{equation}
%	\tilde{f}_m(t) \ge f_{M=2}^{\rm hom}(t;p,q_m^{\rm in}) .
%\end{equation}
%	In addition, by the dominance principle, $f_{M=2}^{\rm hom}(t;p,q_m^{\rm in}) \ge f_{M=2}^{\rm hom}(t;p,q)$.
%	Hence, the result follows. 
%\end{proof}
\subsection{Lower bound}
%In this section we prove that~$\{f_m\}_{m\in {\cal M}}$ and~$f$ are bounded from below from by~$f_{M=2}^{\rm hom}$.

{\bf Proof of Theorem~\ref{thm:f_m>=f_M=2}. }
	To prove the lower bound~\eqref{eq:f_m>=f_M=2} for~$f_m$,
	%Let $[S_m] = 1-f_m$. Then 
	it is sufficient to show that
	\begin{eqnarray*}
		[S_m](t) \leq [S_{M=2}^{\rm hom}](t;p,q) :=1-f_{M=2}^{\rm hom}(t;p,q) 
		=  e^{-pt}\frac{qe^{-pt}-pe^{-qt}}{q-p},
	\end{eqnarray*}
where $[S_m] = 1-f_m$. 
%	By Theorem~\ref{thm:dominance-principle-f_j}, it is sufficient to prove the result in the homogeneous case~\eqref{eq:p_j,q_j-homogeneous}.
%	In that case,	 
%		\begin{equation}
%				\label{eqs:lower-bound-S_M=2-aux}
%			[S_m] = e^{-(p+q)t}+\int_0^te^{-(p+q)\left(t-\tau\right)}\sum_{k \in \cal M}q_{k,m}[S_m,S_k](\tau)\, d\tau ,
%		\end{equation}
%		see~\eqref{eq:S_j-imtegral-identity-upper-bound}.
	By the upper bound in~\eqref{eq:set_e-2pt} and~\eqref{eq:S_j-imtegral-identity-upper-bound}, we have that
		\begin{equation*}
			\begin{aligned}
			 & [S_m]  \leq e^{-(p+q)t}+\int_0^te^{-(p+q)\left(t-\tau\right)}\sum_{k \in \cal M}q_{k,m}e^{-2p\tau}\, d\tau
			= e^{-(p+q)t}+q\int_0^te^{-(p+q)\left(t-\tau\right)}e^{-2p\tau}\, d\tau 
			\\
			&~~~= \left(1-\frac{q}{q-p}\right)e^{-(p+q)t}+\frac{q}{q-p}e^{-2pt} = \left[S_{M=2}^{\rm hom}\right](t;p,q).
		\end{aligned}
	\end{equation*}
	Therefore, we proved~\eqref{eq:f_m>=f_M=2}, which implies~\eqref{eq:f>=ff_M=2}. \qed

%\begin{lemma}
%	  \label{lem:aux-f_m>=f_M=2}
%	Consider the  Bass model~\eqref{eq:Bass-model-heterogeneous-tools} with $p_{j} =p$ for $ j \in \cal M$, $q_m^{\rm in}>0$, and 
%	$q_j = 0$ for $j \in {\cal M} \setminus \{m\}$. 
%	Then 
%	$$
%		{f}_m(t) \ge f_{M=2}^{\rm hom}(t;p,q_m^{\rm in}), \qquad t>0 .
%	$$
%\end{lemma}
%\begin{proof}
%	By construction, for any node~$k$ different from~$m$, there are no incoming edges, and there is a single outgoing edge into~$m$. Therefore,   
%	$m$~is a funnel node of the M-1 subsets $\{A_k\}_{k=1, k \not=m}^M$, where 
%	$A_k:=\{k\}$.
%	Therefore, by the funnel equality for more than 2 sets ({\bf TOMER - need to add this extension to the paper}),
%	$$
%	[S_m](t) = e^{-pt} \prod_{k=1, k \not=m}^M \mathbb{P} \left(X^{A_k}_m(t)=0\right).
%	$$ 
%	Now, $  \mathbb{P}\left(X^{A_k}_m(t)=0\right) = 1-  \mathbb{P}\left(X^{A_k}_m(t)=1\right)$, where
%	$$
%	 \mathbb{P} \left(X^{A_k}_m(t)=1\right)
%	 =
%	   \int_0^t p e^{-p \tau_k} \, \mathbb{P} \left(\text{$m$ adopts
	%	 	in $(\tau_k, t]$} \right)  \, d\tau_k 
% 	=  
% 	\int_0^t p e^{-p \tau_k}\left(1-e^{-q_{k,m}(t-\tau_k)}
%	  \right) \, d\tau_k ,
%	 $$
%	 where $\tau_k$ is the adoption time of~$k$. Explicitly calculating the integral yields {\bf TOMER PLS VERIFY!)}
%	$$
%	\mathbb{P} \left(X^{A_k}_m(t)=0\right) = \frac{q_{k,m} e^{-p t}-p e^{-q_{k,m}t}}{q_{k,m}-p}.
%	.
%	$$ 
%	Therefore, it remains to show that if $q_m^{\rm in} = \sum_{k=1, k \not=m}^M q_{k,m}$ then 
%$$
% e^{-pt} \prod_{k=1, k \not=m}^M \frac{q_{k,m} e^{-p t}-p e^{-q_{k,m}t}}{q_{k,m}-p}  \ge 
% e^{-pt}\frac{q_m^{\rm in} e^{-pt}-pe^{-q_m^{\rm in} t}}{q_m^{\rm in}-p}.
%$$ 
%	
%	
%\end{proof}
%

{\bf Proof of Theoren~\ref{thm:f_j>f_M=2^hom}}
	The only inequality in the proof of Theorem~\ref{thm:f_m>=f_M=2}
	arises from using the  upper bound in~\eqref{eq:set_e-2pt}.
	Therefore,  the lower bound~\eqref{eq:f_m>=f_M=2} for~$[S_m]$ becomes an equality if and only if  $	[S_m,S_k] = e^{-2 pt}$ for all $k \in {\cal M} \setminus m$ for which $q_{k,m}>0$. A minor modification of
	Theorem~\ref{thm:dominance-principle-f_j} shows that  
	$$
	[S_j,S_k] = e^{-2 pt} \quad \iff \quad \text{$j$ and $k$ are not influenced by any other node}.
	$$
	Therefore, \eqref{eq:f_m>f_2nodes}~follows.
%	
%	To prove the claim, note that when~\eqref{eq:p_j,q_j-homogeneous} holds,  
%	the master equation for $[S_j,S_k]$ reads~\citep{fibich2022diffusion}
%	\begin{equation*}
%		\frac{d}{dt}[S_j,S_k]= -2p [S_j,S_k]
%		-\sum_{ l \in {\cal M} \setminus \{k,j\}}  (q_{l,k}+q_{l,j}) [S_{k,j} \cap I_l].
%	\end{equation*}
%	In addition, $[S_{k,j} \cap I_l]>0$ for $t>0$, see~\eqref{eq:0<[S_Omega_cap_In]<1}. Therefore,
%	$\frac{d}{dt}[S_j,S_k]= -2p [S_j,S_k]$ if and only if $\sum_{ l \in {\cal M} \setminus \{k,j\}}  (q_{l,k}+q_{l,j})=0$,
%	which is the case if and only if $j$ and $k$ are not influenced by any other node.\footnote{The claim also follows from Lemma~\ref{lem:strong-dominance-principle-[S_Omega]}.}
	Since $f = \frac{1}{M}\sum_{j=1}^M f_j$ for finite networks and
	$f = \lim_{M\rightarrow \infty }\frac{1}{M}\sum_{j=1}^M f_j$ for infinite networks, 	\eqref{eq:f>f_M=2^hom-thm} also follows. \qed


{\bf Proof of Lemma~\ref{lem:f_M=2-is-tight}. }
	When~$M=2$, this bound is attained by ${\cal N}={\cal N}_{M=2}^{\rm hom}(p,q)$. Moreover, 
	this bound is also attained by any finite or infinite network which is 
	a collection of disjoint pairs of nodes, each of which is of type~${\cal N}_{M=2}^{\rm hom}(p,q)$. \qed

	
\subsection{Asymptotic evaluation of~$\frac{T_{\rm Bass}^{1/2}}{T^{1/2, {\rm hom}}_{M=2}}$}
By~\eqref{eq:f_all-M=2}, $T^{1/2}:=T_{M=2}^{1/2, {\rm hom}}$ is the solution of 
\begin{equation}
\label{eq:t_approx}
e^{-pT^{1/2}}\frac{qe^{-pT^{1/2}}-pe^{-qT^{1/2}}}{q-p} = \frac{1}{2}.
\end{equation}
Let~$X:=e^{-pT^{1/2}}$ and~$\lambda:=\frac{q}{p}$. Then~$e^{-qT^{1/2}}=X^{\frac{q}{p}}=X^{\lambda}$. Plugging this into~\eqref{eq:t_approx}, and noting that~$0<X<1$ and~$\lambda > 0$ gives
\begin{equation*}
\label{eq:X_first}
X^2 - \frac{1}{2} = \frac{p}{q-p}\left(-X^2+X^{\lambda}\right)=O\left(\frac{1}{\lambda}\right), \qquad \lambda\gg 1.
\end{equation*}  
Therefore, 
\begin{equation*}
\label{eq:X_1^2}
X^2 \sim \frac{1}{2}, \qquad \lambda\gg 1.
\end{equation*}
Hence, by the definition of~$X$
\begin{equation*}
\label{eq:T_X}
T^{1/2, {\rm hom}}_{M=2} = \frac{1}{2p}\log(X^{-2})\sim \frac{\log(2)}{2p}, \qquad \lambda\gg 1.
\end{equation*}
Finally, by~\citealp[Lemma 11]{OR-10}, 
$$T_{\rm Bass}^{1/2} = \frac{\log\left(2+\frac{q}{p}\right)}{p+q}\sim \frac{\log\left(\frac{q}{p}\right)}{q},  \qquad \lambda\gg 1,
$$
 and so~\eqref{eq:approx} follows. \qed
%\subsection{``Counter example''}
%In Section~\ref{sec:f_j,f<f_Bass} we showed that the upper bound $f_{\rm Bass}$ holds for any network with $p_j \equiv p$ and~$q_j^{\rm} \equiv q$. In contrast, the lower bound~$f_{\rm 1D}$ does not hold for any network:
%\begin{lemma}
%There exist finite and infinite networks that satisfy $p_j\equiv p$ and $q_j^{\rm in}=q$, for which $f(t;p,q)<f_{\rm 1D}(t;p,q)$ for~$t>0$.
%\end{lemma}
%\begin{proof}
% Consider a pair of nodes, such that the internal effect of each node on its counterpart is $q$. By the indifference principle, the adoption probability of each node is~$f_j^{\rm 1-sided}(t;M=2)$, see~\eqref{eq:fj_one_sided_def}, and by  Lemma~\ref{lem:f_j^one-sided-monotone_in_j}, $f_j^{\rm 1-sided}(t;,M=2) < f_{\rm 1D}(t)$, see~\eqref{eq:f_1_sided_f1D}. This result remain if we take a finite or infinite number of such pairs, each disconnected from all other pairs. Therefore, the result follows.
%\end{proof}
%
%\begin{lemma}
%Consider the  Bass model~\eqref{eq:dbm}. If~$p_j \equiv p$, then~$f_j(t), f(t) > 1-e^{-pt}$.
%\end{lemma}
%\begin{proof}
%The nonadoption probability of an isolated node is~$e^{-pt}$, see~\eqref{eq:e^-pt>S_Bass}. Hence, the result follows from the dominance principle (Theorem~\ref{thm:dominance-principle-f_j}) and~\eqref{eq:nonadoption_prob}. 
%\end{proof}


% Appendix here
% Options are (1) APPENDIX (with or without general title) or
%             (2) APPENDICES (if it has more than one unrelated sections)
% Outcomment the appropriate case if necessary
%
% \begin{APPENDIX}{<Title of the Appendix>}
% \end{APPENDIX}
%
%   or
%
% \begin{APPENDICES}
% \section{<Title of Section A>}
% \section{<Title of Section B>}
% etc
% \end{APPENDICES}

%%




\section*{Acknowledgements}
We thank Eilon Solan for useful comments.
 

%% The Appendices part is started with the command \appendix;
%% appendix sections are then done as normal sections
%% \appendix

%% \section{}
%% \label{}

%% If you have bibdatabase file and want bibtex to generate the
%% bibitems, please use
%%
  \bibliographystyle{apalike} 
  \bibliography{diffusion.bib}



%% else use the following coding to input the bibitems directly in the
%% TeX file.



\end{document}
\endinput
%%
%% End of file `elsarticle-template-num.tex'.

\begin{thebibliography}{}

\bibitem[Albert et~al., 2000]{Albert-00}
Albert, R., Jeong, H., and Barab\'asi, A. (2000).
\newblock Error and attack tolerance of complex networks.
\newblock {\em Nature}, 406:378--382.

\bibitem[Anderson and May, 1992]{Anderson-92}
Anderson, R. and May, R. (1992).
\newblock {\em Infectious Diseases of Humans}.
\newblock Oxford University Press, Oxford.

\bibitem[Barab{\'a}si and Albert, 1999]{barabasi1999emergence}
Barab{\'a}si, A. and Albert, R. (1999).
\newblock Emergence of scaling in random networks.
\newblock {\em {S}cience}, 286:509--512.

\bibitem[Bass, 1969]{Bass-69}
Bass, F. (1969).
\newblock A new product growth model for consumer durables.
\newblock {\em Management science}, 15:1215--1227.

\bibitem[{Bollinger} and {Gillingham}, 2012]{Bollinger-12}
{Bollinger}, B. and {Gillingham}, K. (2012).
\newblock Peer effects in the diffusion of solar photovoltaic panels.
\newblock {\em Marketing Science}, 31:900--912.

\bibitem[Delre et~al., 2007]{Delre-07}
Delre, S., Jager, W., Bijmolt, T., and Janssen, M. (2007).
\newblock Targeting and timing promotional activities: {A}n agent-based model
  for the takeoff of new products.
\newblock {\em Journal of business research}, 60:826--835.

\bibitem[Erd{\H{o}}s and R{\'e}nyi, 1960]{erdHos1960evolution}
Erd{\H{o}}s, P. and R{\'e}nyi, A. (1960).
\newblock On the evolution of random graphs.
\newblock {\em Publications of the Mathematical Institute of the Hungarian
  Academy of Sciences}, 5:17--60.

\bibitem[Fibich et~al., 2012]{PNAS-12}
Fibich, G., Gavious, A., and Solan, E. (2012).
\newblock Averaging principle for second-order approximation of heterogeneous
  models with homogeneous models.
\newblock {\em Proceedings of the National Academy of Sciences},
  109:19545--19550.

\bibitem[Fibich and Gibori, 2010]{OR-10}
Fibich, G. and Gibori, R. (2010).
\newblock Aggregate diffusion dynamics in agent-based models with a spatial
  structure.
\newblock {\em Operations Research}, 58:1450--1468.

\bibitem[Fibich and Golan, 2022]{fibich2022diffusion}
Fibich, G. and Golan, A. (2022).
\newblock Diffusion of new products with heterogeneous consumers.
\newblock {\em Mathematics of Operations Research}.

\bibitem[Fibich et~al., 2023]{Bass-compartmental-limits-22}
Fibich, G., Golan, A., and Schochet, S. (2023).
\newblock Compartmental limit of discrete {B}ass models on networks.
\newblock {\em Discrete and Continuous Dynamical Systems - B}, 28:3052--3078.

\bibitem[Fibich et~al., rint]{Bass-monotone-convergence-23}
Fibich, G., Golan, A., and Schochet, S. (Preprint).
\newblock Monotone convergence of discrete {B}ass models.

\bibitem[Fibich and Levin, rint]{Funnel2022}
Fibich, G. and Levin, T. (Preprint).
\newblock Funnel theorems for spreading on networks.
\newblock {\em available at \url{http://arxiv.org/abs/2308.13034}}.

\bibitem[Fibich et~al., 2019]{Bass-boundary-18}
Fibich, G., Levin, T., and Yakir, O. (2019).
\newblock Boundary effects in the discrete bass model.
\newblock {\em SIAM Journal on Applied Mathematics}, 79:914--737.

\bibitem[Garber et~al., 2004]{Garber-04}
Garber, T., Goldenberg, J., Libai, B., and Muller, E. (2004).
\newblock From density to destiny: Using spatial dimension of sales data for
  early prediction of new product success.
\newblock {\em Marketing Science}, 23:419--428.

\bibitem[Garcia, 2005]{Garcia-05}
Garcia, R. (2005).
\newblock Uses of agent-based modeling in innovation/new product development
  research.
\newblock {\em Journal of Product Innovation Management}, 22:380--398.

\bibitem[Goldenberg et~al., 2001]{GLM-01}
Goldenberg, J., Libai, B., and Muller, E. (2001).
\newblock Using complex systems analysis to advance marketing theory
  development. (special issue on emergent and co-evolutionary processes in
  marketing.).
\newblock {\em Acad. Market. Sci. Rev.}, 9:1--19.

\bibitem[Goldenberg et~al., 2010]{Goldenberg-10}
Goldenberg, J., Libai, B., and Muller, E. (2010).
\newblock The chilling effects of network externalities.
\newblock {\em International Journal of Research in Marketing}, 27:4--15.

\bibitem[Goldenberg et~al., 2008]{Goldenberg-08}
Goldenberg, J., Lowengart, O., and Shapira, D. (2008).
\newblock Zooming in: Self-emergence of movements in new product growth.
\newblock {\em Marketing Science}, 28:274--292.

\bibitem[Graziano and Gillingham, 2015]{Graziano-15}
Graziano, M. and Gillingham, K. (2015).
\newblock Spatial patterns of solar photovoltaic system adoption: {T}he
  influence of neighbors and the built environment.
\newblock {\em Journal of Economic Geography}, 15:815--839.

\bibitem[Hopp, 2004]{Hopp-04}
Hopp, W. (2004).
\newblock Ten most influential papers of management science's first fifty
  years.
\newblock {\em Management Science}, 50:1763--1764.

\bibitem[Jackson, 2008]{Jackson-08}
Jackson, M. (2008).
\newblock {\em Social and Economic Networks}.
\newblock Princeton University Press, Princeton and Oxford.

\bibitem[Kiss et~al., 2017]{kiss2017mathematics}
Kiss, I., Miller, J., and Simon, P. (2017).
\newblock {\em Mathematics of epidemics on networks}.
\newblock Springer.

\bibitem[Mahajan et~al., 1993]{Mahajan-93}
Mahajan, V., Muller, E., and Bass, F. (1993).
\newblock New-product diffusion models.
\newblock {\em Handbooks in operations research and management science},
  5:349--408.

\bibitem[Muller and Peres, 2019]{Muller-Peres-19}
Muller, E. and Peres, R. (2019).
\newblock The effect of social networks structure on innovation performance: A
  review and directions for research.
\newblock {\em International Journal of Research in Marketing}, 36:3--19.

\bibitem[Niu, 2002]{Niu-02}
Niu, S. (2002).
\newblock A stochastic formulation of the {B}ass model of new product
  diffusion.
\newblock {\em Mathematical problems in Engineering}, 8:249--263.

\bibitem[Pastor-Satorras and Vespignani, 2001]{Pastor-Satorras-01}
Pastor-Satorras, R. and Vespignani, A. (2001).
\newblock Epidemic spreading in scale-free networks.
\newblock {\em Physical review letters}, 86:3200--3203.

\bibitem[Peres, 2014]{Peres-14}
Peres, R. (2014).
\newblock The impact of network characteristics on the diffusion of
  innovations.
\newblock {\em Physica A}, 402:330--343.

\bibitem[Rand and Rust, 2011]{Rand-11}
Rand, W. and Rust, R. (2011).
\newblock Agent-based modeling in marketing: Guidelines for rigor.
\newblock {\em International Journal of research in Marketing}, 28:181--193.

\bibitem[Scellato et~al., 2011]{scellato2011socio}
Scellato, S., Noulas, A., Lambiotte, R., and Mascolo, C. (2011).
\newblock Socio-spatial properties of online location-based social networks.
\newblock In {\em Proceedings of the International AAAI Conference on Web and
  Social Media}, volume~5, pages 329--336.

\bibitem[Strang and Soule, 1998]{Strang-98}
Strang, D. and Soule, S. (1998).
\newblock Diffusion in organizations and social movements: From hybrid corn to
  poison pills.
\newblock {\em Annual review of sociology}, 24:265--290.

\bibitem[Watts and Strogatz, 1998]{watts1998collective}
Watts, D. and Strogatz, S. (1998).
\newblock Collective dynamics of ‘small-world’ networks.
\newblock {\em {N}ature}, 393:440--442.

\end{thebibliography}
